\chapter{Interpretable Regression with Checked Physical Predictions}
\label{ch:interpretable}

\section{Rationale}

An activated sludge surrogate used for process analysis should do more than reproduce treatment outputs accurately. It should also allow an engineer to examine how operating settings and influent composition contribute to those predictions. An error score alone cannot explain why a predicted nutrient concentration changes when aeration or loading changes, nor can it reveal whether the underlying component response remains physically meaningful. The prediction must retain the substrate, nutrient, biomass, and solids components required for both treatment reporting and material accounting \citep{Henze2006}. The research problem is therefore to obtain an inspectable approximation of the complete activated sludge state whose returned concentrations also satisfy the adopted physical requirements.

Directly inspectable models provide a basis for scrutinizing predictions used in consequential decisions \citep{Rudin2019}. Polynomial response-surface methods offer a treatment-specific precedent for representing operating effects and interactions explicitly \citep{Nazlabadi2021}. The nonlinear libraries of \citet{Brunton2016} and the combined process and learning models reviewed by \citet{Shen2023} likewise show how mathematical structure can be connected with process knowledge. These approaches establish the value of explicit representations, but they do not by themselves provide an interpretable prediction of the complete activated sludge component state after physical enforcement. The first gap therefore concerns how operating and loading effects can remain identifiable throughout the full prediction calculation.

This distinction between operation and loading is especially important for nutrient removal. Carbon availability affects denitrification and phosphorus storage, while oxygen availability affects carbon oxidation and nitrification. \citet{Wentzel1992} describe these interconnected pathways, and \citet{Guerrero2011} examine competition for carbon during nitrogen and phosphorus removal. Consequently, the predicted effect of increasing aeration may depend on the incoming ammonium concentration, organic substrate, or both. Separate operating, loading, curvature, and interaction terms make these relationships visible within the fitted model. Their interpretation, however, remains an interpretation of the statistical response and does not establish a unique biological mechanism.

Relationships among effluent components introduce a second layer of interpretation. A term associated with influent substrate can influence several predicted concentrations through learned component coupling. Its coefficient in the regression driver is therefore not necessarily the final sensitivity of substrate, biomass, or another returned component. Both the units of the coefficient and its position within the calculation matter. The structured treatment model of \citet{Selerio2026} provides the basis for separating the explicit regression driver from the coupled raw effluent state. Examining these stages separately makes the fitted associations traceable without interpreting coupling coefficients as kinetic parameters.

The second gap concerns what happens to interpretation after physical enforcement. Projection changes the raw component state before COD, TN, TP, and TSS are calculated. A concentration bound may be inactive at one operating point and active at another, so the local behavior of the returned prediction can differ from the slope suggested by the raw regression. Interpretation must therefore follow the complete sequence from the driver, through component coupling, to the projected effluent \citep{Selerio2026}. A visible polynomial is not sufficient if the explanation stops before the operation that produces the state ultimately used by the engineer.

Physical requirements must also be distinguished from fitting preferences. Penalties can discourage balance violations or negative concentrations during training, but they do not guarantee that every new prediction satisfies both requirements \citep{Hansen2024}. An explicit polynomial can still return negative components even when all training states are non-negative. Mass conservation and non-negativity therefore require a defined enforcement procedure followed by checks on the final returned state. The joint projection of \citet{KircherVotsmeier2025} demonstrates the importance of treating these conditions together. For an inspectable activated sludge regression, the enforcement stage must therefore be both physically checked and included in the explanation of the prediction.

The practical value of this structure also depends on the amount of available training data. A compact representation of operating effects may be attractive when mechanistic calculations limit the number of training cases, but it need not outperform a more flexible learner when a larger simulation design is available. Learning curves make this distinction measurable \citep{VieringLoog2023}. Held-out prediction error, fitting effort, deployment effort, coefficient structure, and physical checks are therefore assessed together. The purpose is not to assume that a transparent equation is automatically the best approximation, but to determine when inspectability is accompanied by useful predictive performance.

The study presents invariant-constrained second-order regression (ICSOR), which combines an explicit polynomial driver, learned component coupling, and checked prediction \citep{Selerio2026}. Deployment here means evaluating a fitted model at a new influent and operating condition and returning its physically checked component state. The study addresses the third research question in Section~\ref{sec:research_questions} by examining the complete response, the fitted terms that generate it, and its sensitivities before and after projection. Its contribution is an inspectable component-state formulation with a traceable physical enforcement procedure. The assessment identifies the data conditions under which this structure is useful and the limits that remain when its coefficients are interpreted. It does not presume causal coefficient meanings or superior predictive accuracy at every training size.

\section{Theory}

The prediction target is the effluent component vector \(c_{out}\in\mathbb{R}^{F}\). The adopted process basis contains \(F=20\) components and \(P=28\) reactions. The operating vector \(u\in\mathbb{R}^{M_{op}}\) has \(M_{op}=2\) entries representing HRT and aeration level, while the full influent state is \(c_{in}\in\mathbb{R}^{F}\). Component concentrations retain their physical units throughout the calculation. The four reported water-quality composites are obtained from a fixed reporting map \(I_{\mathrm{comp}}\in\mathbb{R}^{4\times F}\), which is applied only after the component state has been predicted.

\subsection{Mass conservation in the adopted state space}

The connection between a component balance and a regression constraint can first be illustrated with a hypothetical conversion of one component into another. Suppose one unit of component 1 is converted into one unit of component 2, with both expressed in a common inventory unit. Their changes satisfy
\begin{equation}
\equationname{Illustrative component changes for a two-component reaction}
\label{eq:icsor_scalar_inventory_example}
c_{out,1}-c_{in,1}=-\xi,
\qquad
c_{out,2}-c_{in,2}=\xi.
\end{equation}
Adding the two equations eliminates the unknown conversion amount and gives the relation \(c_{out,1}+c_{out,2}=c_{in,1}+c_{in,2}\). For an illustrative influent \((6,4)^\top\) and conversion amount 2, the effluent is \((4,6)^\top\), and both states have a total inventory of 10. Many other non-negative pairs can satisfy the same total. The balance therefore defines a set of admissible regression outputs rather than a unique state.

The same calculation can be written using the reduced matrices
\begin{equation}
\equationname{Illustrative stoichiometric matrix and mass conservation operator}
\label{eq:icsor_small_inventory_matrices}
\nu=\begin{bmatrix}-1&1\end{bmatrix},
\qquad
A=\begin{bmatrix}1&1\end{bmatrix},
\qquad
A\nu^\top=1(-1)+1(1)=0.
\end{equation}
These hypothetical matrices use only two components to make the operation visible. The adopted reactor applies the same principle to the full 20-component basis, reaction matrix, and operating domain.

The stoichiometric matrix \(\nu\in\mathbb{R}^{P\times F}\) stores reactions as rows and components as columns. Let \(V_0\in\mathbb{R}^{F\times K}\) span its right null space, where \(K=5\). The invariant operator is
\begin{equation}
\equationname{Structured-regression mass conservation operator from the null space}
\label{eq:icsor_invariant_operator}
A=V_0^\top\in\mathbb{R}^{K\times F},
\qquad \nu V_0=0,
\qquad A\nu^\top=0.
\end{equation}
Because the columns of \(V_0\) are independent, \(A\) has full row rank. When the null-space basis is orthonormal, \(AA^\top=I_K\). Another full-rank basis would describe the same mass conservation conditions, although its row scaling would change the numerical magnitude of a soft invariant penalty.

For the control volume and boundary accounting represented by the adopted process model, a reaction-induced component change can be written as
\begin{equation}
\equationname{Effluent change in stoichiometric process coordinates}
\label{eq:icsor_stoichiometric_change}
c_{out}-c_{in}=\nu^\top\xi,
\qquad \xi\in\mathbb{R}^{P}.
\end{equation}
The reaction-progress vector \(\xi\) represents the net contribution of the modeled reactions in concentration-equivalent units. Multiplication by \(A\) removes this unknown reaction contribution and yields
\begin{equation}
\equationname{Mass conservation equality for the structured-regression effluent}
\label{eq:icsor_invariant_balance}
A(c_{out}-c_{in})=0.
\end{equation}
This construction follows the mass conservation logic used in the photochemical surrogate framework of \citet{SturmWexler2020}. \citet{SturmWexler2022} further show how a stoichiometric transformation can connect learned reaction information with state changes that satisfy mass conservation. The same linear algebra is used here, but the conserved combinations are defined by the activated sludge reaction network.

Equation~\eqref{eq:icsor_invariant_balance} applies to the modeled boundary in Equation~\eqref{eq:icsor_stoichiometric_change}. External gas transfer, chemical dosing, solids withdrawal, and other boundary fluxes must enter the balance where they occur. If an explicit boundary contribution \(s\) is present, the corresponding relation becomes \(A(c_{out}-c_{in}-s)=0\).

For a new influent state, checked deployment therefore uses the feasible component set
\begin{equation}
\equationname{Structured-regression feasible set for mass conservation and non-negativity}
\label{eq:icsor_feasible_set}
\mathcal{F}(c_{in})=
\left\{c\in\mathbb{R}^{F}\mid Ac=Ac_{in},\ c\geq\mathbf{0}\right\}.
\end{equation}
If \(c_{in}\geq\mathbf{0}\), the set contains \(c_{in}\) itself and is therefore nonempty. Its members satisfy both the modeled linear invariants and component non-negativity within the declared system boundary.

The reported water-quality composites provide another reason to perform these checks in component space. Distinct component states can produce the same COD, TN, TP, and TSS values, and a negative component can be hidden by positive contributions to an aggregate. When \(I_{\mathrm{comp}}\) is entrywise non-negative, a non-negative component state necessarily produces non-negative reported composites, but the reverse implication does not generally hold. Mass conservation and non-negativity are therefore assessed before the state is collapsed into the four reported quantities \citep{Henze2006,Selerio2026}.

\subsection{A partitioned second-order driver}

\subsubsection{Separating operation from loading}

Operating conditions and influent composition have different engineering roles. HRT and aeration determine aspects of the reactor environment, whereas influent components determine the material entering that environment. This distinction is consistent with the operating-sensitivity analysis of \citet{Araujo2013} and with the separation of operational and influent uncertainty examined by \citet{Machado2014}.

Consider first two hypothetical operating entries \(u_1,u_2\) and two influent entries \(c_1,c_2\). A second-order model includes the individual entries, their squares, and their pairwise products. The operating products are \(u_1^2,u_1u_2,u_2u_1,u_2^2\). Both cross-term orders are retained as an accounting convention even though \(u_1u_2=u_2u_1\).

For two column vectors, the Kronecker product multiplies the complete second vector by each entry of the first vector in sequence. Thus
\begin{equation}
\equationname{Illustrative operating, influent, and mixed Kronecker products}
\label{eq:icsor_small_kronecker}
\begin{aligned}
u\otimes u&=\begin{bmatrix}u_1u_1\\u_1u_2\\u_2u_1\\u_2u_2\end{bmatrix},
&c\otimes c&=\begin{bmatrix}c_1c_1\\c_1c_2\\c_2c_1\\c_2c_2\end{bmatrix},\\
u\otimes c&=\begin{bmatrix}u_1c_1\\u_1c_2\\u_2c_1\\u_2c_2\end{bmatrix}.
\end{aligned}
\end{equation}
For the illustrative values \(u=(2,3)^\top\) and \(c=(4,5)^\top\), these three blocks become \((4,6,6,9)^\top\), \((16,20,20,25)^\top\), and \((8,10,12,15)^\top\), respectively. Including the baseline and first-order entries gives 17 ordered features in this reduced example. The adopted model uses the same construction for the full input vector.

The complete feature vector is
\begin{equation}
\equationname{Second-order driver feature vector and its dimension}
\label{eq:icsor_features}
\phi(u,c_{in})=
\begin{bmatrix}
1\\u\\c_{in}\\u\otimes u\\c_{in}\otimes c_{in}\\u\otimes c_{in}
\end{bmatrix}
\in\mathbb{R}^{D},
\qquad
D=1+M_{op}+F+M_{op}^{2}+F^{2}+M_{op}F.
\end{equation}
The symbol \(\otimes\) denotes the Kronecker product. For the adopted dimensions, \(D=467\). The ordered expansion includes both \(u_i u_j\) and \(u_j u_i\), as well as both \(c_{in,i}c_{in,j}\) and \(c_{in,j}c_{in,i}\). Although these paired features have identical numerical values, their coefficient representation requires the symmetry convention introduced below.

The multiplication of features by regression coefficients can be read one entry at a time. For an illustrative first driver, take a baseline coefficient of 1, operating coefficients \((2,-1)\), and influent coefficients \((0.5,0)\). Let the operating quadratic row be \((0.1,0.1,0.1,0)\), the influent quadratic row be \((0,0,0,-0.05)\), and the operating-load row be \((0,0,0.25,0)\). Applying these coefficients to the numerical feature blocks gives
\begin{equation}
\equationname{Numerical evaluation of an illustrative regression driver}
\label{eq:icsor_driver_multiplication_example}
\begin{aligned}
r_1={}&1+\begin{bmatrix}2&-1\end{bmatrix}\begin{bmatrix}2\\3\end{bmatrix}
+\begin{bmatrix}0.5&0\end{bmatrix}\begin{bmatrix}4\\5\end{bmatrix}\\
&+\begin{bmatrix}0.1&0.1&0.1&0\end{bmatrix}\begin{bmatrix}4\\6\\6\\9\end{bmatrix}\\
&+\begin{bmatrix}0&0&0&-0.05\end{bmatrix}\begin{bmatrix}16\\20\\20\\25\end{bmatrix}\\
&+\begin{bmatrix}0&0&0.25&0\end{bmatrix}\begin{bmatrix}8\\10\\12\\15\end{bmatrix}\\
={}&1+1+2+1.6-1.25+3=7.35.
\end{aligned}
\end{equation}
Each row-vector product is simply a sum of coefficient--feature products. The calculation is therefore equivalent to an ordinary linear regression evaluated on an expanded nonlinear feature library. A second hypothetical driver \(r_2=-u_1+u_2+0.2c_2\) would equal 2 at the same inputs. The subsequent coupled relation converts these driver values into predicted component concentrations.

Before grouping the coefficients into matrices, the same first driver can be written in scalar polynomial form as
\begin{equation}
\equationname{Scalar polynomial form of the illustrative regression driver}
\label{eq:icsor_scalar_driver_example}
r_1=1+2u_1-u_2+0.5c_1+0.1u_1^2
+0.2u_1u_2-0.05c_2^2+0.25u_2c_1.
\end{equation}
The coefficient 0.2 of \(u_1u_2\) is the sum of the two ordered coefficients of 0.1. More generally, the driver for component \(f\) is \(r_f=\sum_{\ell=1}^{D}B_{f\ell}\phi_\ell\). Evaluating this sum for all components gives
\begin{equation}
\equationname{Matrix mapping from second-order features to the regression driver}
\label{eq:icsor_driver}
r(u,c_{in})=B\phi(u,c_{in}),
\qquad B\in\mathbb{R}^{F\times D}.
\end{equation}
The driver is an intermediate statistical quantity supplied to the coupled component relation. Its coefficients are unrestricted in sign because the fitted response may increase or decrease with an input. Non-negativity is imposed later on the component state, where that requirement has a physical meaning.

The same driver can be partitioned into interpretable feature groups:
\begin{equation}
\equationname{Driver decomposition into operating, influent, and interaction blocks}
\label{eq:icsor_driver_blocks}
\begin{aligned}
r(u,c_{in})={}&b+W_u u+W_{in}c_{in}
+\Theta_{uu}(u\otimes u)\\
&+\Theta_{cc}(c_{in}\otimes c_{in})
+\Theta_{uc}(u\otimes c_{in}).
\end{aligned}
\end{equation}
Here \(b\in\mathbb{R}^{F}\) is the baseline term. The matrices \(W_u\in\mathbb{R}^{F\times M_{op}}\) and \(W_{in}\in\mathbb{R}^{F\times F}\) contain the first-order operating and influent terms. The matrices \(\Theta_{uu}\in\mathbb{R}^{F\times M_{op}^{2}}\) and \(\Theta_{cc}\in\mathbb{R}^{F\times F^{2}}\) represent curvature within the two input groups, while \(\Theta_{uc}\in\mathbb{R}^{F\times M_{op}F}\) contains operating-load interactions. This partition makes the statistical roles of operation, loading, curvature, and interaction explicit.

\subsubsection{Illustrative operating and loading effects}

Consider an illustrative driver for one component with aeration level \(a\) and influent ammonium concentration \(n\):
\begin{equation}
\equationname{Illustrative driver with aeration, nitrogen, and interaction effects}
\label{eq:icsor_example_driver}
r_f=\beta_0+\beta_a a+\beta_n n+\beta_{aa}a^2+\beta_{an}an.
\end{equation}
Its derivative with respect to aeration is \(\partial r_f/\partial a=\beta_a+2\beta_{aa}a+\beta_{an}n\). The modeled aeration response therefore depends not only on the linear aeration coefficient, but also on the current aeration setting and the ammonium loading. The interaction term makes the local aeration effect conditional on the influent condition.

The same reasoning extends to other simultaneous loads. A carbon-ammonium interaction allows the fitted association with carbon availability to depend on nitrogen loading. A phosphorus-storage interaction allows the fitted response to dissolved phosphate to depend on a particulate storage pool. These terms provide a transparent statistical approximation of nonlinear treatment behavior involving the interconnected microbial pathways described by \citet{Henze1999}.

The second-order library therefore represents smooth local curvature and pairwise interactions while keeping those terms directly inspectable. Explicit nonlinear libraries provide a general basis for examining such structures \citep{Brunton2016}, while the response-surface applications reviewed by \citet{Nazlabadi2021} provide a wastewater-treatment precedent for using polynomial operating terms and interactions.

\subsubsection{Symmetry and the meaning of cross terms}

The ordered quadratic expansion contains mirrored cross terms, but those pairs are not separately identifiable from the polynomial response. For a two-entry vector,
\begin{equation}
\equationname{Expansion of a two-variable quadratic form}
\label{eq:icsor_scalar_quadratic}
\begin{bmatrix}x_1&x_2\end{bmatrix}
\begin{bmatrix}q_{11}&q_{12}\\q_{21}&q_{22}\end{bmatrix}
\begin{bmatrix}x_1\\x_2\end{bmatrix}
=q_{11}x_1^2+(q_{12}+q_{21})x_1x_2+q_{22}x_2^2.
\end{equation}
Only the sum of the two off-diagonal entries affects the polynomial. Replacing them by their average therefore leaves the prediction unchanged. For example,
\begin{equation}
\equationname{Numerical example of quadratic-coefficient symmetrization}
\label{eq:icsor_symmetry_numbers}
Q=\begin{bmatrix}0.1&0.3\\-0.1&0\end{bmatrix},
\qquad
\frac{Q+Q^\top}{2}=\begin{bmatrix}0.1&0.1\\0.1&0\end{bmatrix}.
\end{equation}
At \(x=(2,3)^\top\), the first matrix gives \(0.4+1.8-0.6=1.6\), while the symmetric matrix gives \(0.4+0.6+0.6=1.6\). The two matrices therefore make the same prediction even though their individual off-diagonal coefficients differ.

Their squared coefficient sums are not the same. The original matrix has squared sum \(0.01+0.09+0.01=0.11\), while the symmetric part has squared sum \(0.01+0.01+0.01=0.03\). The removed skew-symmetric part contributes 0.08 to the coefficient norm without contributing anything to the prediction. A positive coefficient penalty therefore favors the symmetric representation.

For component \(f\), the corresponding rows of \(\Theta_{uu}\) and \(\Theta_{cc}\) can be reshaped into square matrices \(Q_{uu}^{(f)}\) and \(Q_{cc}^{(f)}\). Any quadratic term can then be written as \(x^\top Qx\), and for any real square matrix,
\begin{equation}
\equationname{Symmetric coefficient matrix of a quadratic form}
\label{eq:icsor_quadratic_symmetry}
x^\top Qx=x^\top\operatorname{sym}(Q)x,
\qquad
\operatorname{sym}(Q)=\frac{Q+Q^\top}{2}.
\end{equation}
The skew-symmetric part contributes nothing to the predicted value but can increase the coefficient norm because
\begin{equation}
\equationname{Frobenius-norm decomposition into symmetric and skew parts}
\label{eq:icsor_quadratic_norm}
\|Q\|_F^2=
\|\operatorname{sym}(Q)\|_F^2+
\|\operatorname{skew}(Q)\|_F^2.
\end{equation}
Replacing each quadratic slice by its symmetric part therefore preserves all fitted values and cannot increase a positive ridge penalty. The projection is applied after each driver update. With a strictly positive driver ridge weight, the exact minimum-norm solution contains no skew-symmetric contribution.

This symmetry convention also determines how coefficient displays should be interpreted. If \(q_{ij}=q_{ji}\), the two displayed off-diagonal cells jointly contribute \(2q_{ij}x_i x_j\). Each displayed cell therefore represents half of the coefficient of the corresponding unordered interaction term. For example, if \(q_{12}=q_{21}=0.3\), the polynomial contains \(0.6x_1x_2\). Reading either cell alone as the complete interaction coefficient would understate the cross-term contribution by half.

Each component therefore has \(M_{op}(M_{op}+1)/2=3\) distinct operating quadratic terms and \(F(F+1)/2=210\) distinct influent quadratic terms. Together with the baseline, first-order terms, and operating-load interactions, the model contains 276 distinct polynomial terms per component. These counts define the structural dimension of the interpretable feature library.

\subsection{Learned coupling among effluent components}

The driver does not predict each effluent component independently. Instead, the model allows the component channels to interact through a learned simultaneous relation. Consider two hypothetical component equations with coupling coefficient 0.2 in both directions:
\begin{equation}
\equationname{Illustrative coupled equations for two effluent components}
\label{eq:icsor_scalar_coupling_example}
c_1=4+0.2c_2,
\qquad c_2=2+0.2c_1.
\end{equation}
Substituting the second equation into the first gives \(c_1=4+0.2(2+0.2c_1)\), or \(0.96c_1=4.4\). Thus \(c_1=55/12\) and \(c_2=35/12\), approximately \((4.5833,2.9167)^\top\). These values illustrate the arithmetic of a simultaneous statistical relation: each predicted component depends on both the driver and the other component.

Moving the coupled terms to the left gives the equivalent matrix system
\begin{equation}
\equationname{Illustrative coupling matrix and coupled-state system}
\label{eq:icsor_small_coupling_matrices}
\Gamma=\begin{bmatrix}0&0.2\\0.2&0\end{bmatrix},
\qquad
R=\begin{bmatrix}1&-0.2\\-0.2&1\end{bmatrix},
\qquad
R\begin{bmatrix}c_1\\c_2\end{bmatrix}=\begin{bmatrix}4\\2\end{bmatrix}.
\end{equation}
The determinant of \(R\) is \(1-0.2^2=0.96\), and
\begin{equation}
\equationname{Inverse of the illustrative two-component coupling system}
\label{eq:icsor_small_coupling_inverse}
R^{-1}=\frac{1}{0.96}\begin{bmatrix}1&0.2\\0.2&1\end{bmatrix}.
\end{equation}
Multiplication by this inverse gives the same solution as the scalar substitution. More importantly for interpretation, the inverse shows how a change in one driver can propagate into several predicted component channels.

The full regression uses a coupling matrix \(\Gamma\in\mathbb{R}^{F\times F}\). Define
\begin{equation}
\equationname{General coupled-state relation and system matrix}
\label{eq:icsor_coupled_relation}
R=I_F-\Gamma,
\qquad Rc=r(u,c_{in}).
\end{equation}
Equivalently, the equation for component \(f\) is \(c_f=r_f+\sum_{j\neq f}\Gamma_{fj}c_j\). This formulation separates the part of the response driven directly by the inputs from the dependence learned among the predicted component channels. The diagonal of \(\Gamma\) is fixed to zero so that self-dependence is not represented redundantly.

The sign of \(\Gamma_{fj}\) describes reinforcement or compensation within the fitted simultaneous relation. Activated sludge reactions involve several components simultaneously, and the stoichiometric matrix specifies those process transformations \citep{Henze2006}. The coupling matrix, however, is learned from state patterns rather than imposed from kinetics. Its entries therefore describe fitted dependence arising from correlated transformations, loading patterns, and the allocation of effects between the driver and coupling blocks \citep{Selerio2026}.

When \(R\) is nonsingular, the raw component prediction is \(R^{-1}r\). This multiplication can redistribute or amplify the response encoded by the driver. Consequently, interpretation requires examination of the driver, the coupling matrix, and their combined raw mapping rather than any one of these objects in isolation.

The coupling matrix is subject to
\begin{equation}
\equationname{Diagonal, magnitude, and conditioning safeguards for component coupling}
\label{eq:icsor_coupling_safeguards}
\Gamma_{ff}=0,
\qquad |\Gamma_{fj}|\leq\gamma_{abs}\quad(f\neq j),
\qquad \kappa_2(I_F-\Gamma)\leq\kappa_{max}.
\end{equation}
The condition number is the ratio of the largest to the smallest singular value of the coupled-system matrix. A nearly singular system can amplify small errors in the driver or in numerical computation \citep{BuzziFerraris2013}. The condition-number check therefore controls the behavior of the complete coupling system independently of the bounds imposed on individual entries.

The small symmetric matrix in Equation~\eqref{eq:icsor_small_coupling_matrices} makes this distinction visible. Along the direction \((1,1)^\top\), multiplication by \(R\) scales the vector by 0.8; along \((1,-1)^\top\), it scales the vector by 1.2. These orthogonal directions give singular values 0.8 and 1.2, so the condition number is \(1.2/0.8=1.5\).

If the two illustrative coupling coefficients were instead 0.99, the same directions would have scales 0.01 and 1.99, giving a condition number of 199. A driver error of magnitude \(\epsilon\) in each component along the first direction would then produce a state error of \(100\epsilon\) in each component. This hypothetical example shows how strong coupling can greatly amplify a small perturbation.

Conversely, small individual entries do not guarantee that the complete matrix is well conditioned. In a hypothetical four-component system with every off-diagonal coupling equal to \(1/3\), each row sum of \(\Gamma\) equals one. Thus \(\Gamma\mathbf{1}=\mathbf{1}\) and \((I-\Gamma)\mathbf{1}=0\). Every coupling entry is smaller than the adopted bound of 0.368, yet the resulting system matrix is singular. The entry bound and condition-number check therefore protect against different failure modes.

Shrinking a proposed coupling matrix toward zero provides a natural safe direction because \(I_F-0=I_F\), whose condition number is one. A prescribed \(\kappa_{max}>1\) therefore admits the zero matrix as a fallback. The condition-number requirement defines a generally nonconvex admissible set, so it is evaluated after the convex box-constrained coupling update proposes a candidate.

\section{Methodology}

\subsection{Fitting in prediction space}

\subsubsection{The coupled training objective}

The fitting procedure estimates coefficients that approximate the mechanistic effluent concentrations while retaining the structured relation described above. It also introduces auxiliary fitted states, which are concentration variables optimized for the training cases. These auxiliary states make it possible to express prediction error, physical guidance, and coupled-system consistency within one objective. They are fitting variables and should not be confused with the component predictions returned for new inputs after the model has been trained.

For any matrix \(M\) with entries \(M_{ij}\), the squared Frobenius norm is
\begin{equation}
\equationname{Squared Frobenius norm as a sum of coefficient squares}
\label{eq:icsor_frobenius_definition}
\|M\|_F^2=\sum_i\sum_j M_{ij}^2.
\end{equation}
For example, a residual matrix \(\begin{bmatrix}1&-2\\3&0\end{bmatrix}\) has squared norm \(1^2+(-2)^2+3^2+0^2=14\). The corresponding norm is \(\sqrt{14}\), whereas a least-squares objective uses the squared value 14. Squaring the residuals prevents positive and negative errors from canceling when they are aggregated.

Let \(i\) index observations, \(f\) components, \(k\) invariants, and \(\ell\) features. The fitted concentration is \(\widehat c_{if}\), and \(\phi_{i\ell}\) is the corresponding feature value. Written explicitly, the five contributions to the training criterion are
\begin{equation}
\equationname{Scalar expansion of the structured-regression training objective}
\label{eq:icsor_training_scalar}
\begin{aligned}
J={}&\sum_{i=1}^{N}\sum_{f=1}^{F}
\left(c_{out,if}-\widehat c_{if}\right)^2\\
&+\lambda_{inv}\sum_{i=1}^{N}\sum_{k=1}^{K}
\left[\sum_{f=1}^{F}A_{kf}(\widehat c_{if}-c_{in,if})\right]^2\\
&+\lambda_{sys}\sum_{i=1}^{N}\sum_{f=1}^{F}
\left[\widehat c_{if}-\sum_{j=1}^{F}\Gamma_{fj}\widehat c_{ij}
-\sum_{\ell=1}^{D}B_{f\ell}\phi_{i\ell}\right]^2\\
&+\lambda_B\sum_{f=1}^{F}\sum_{\ell=1}^{D}B_{f\ell}^2
+\lambda_{\Gamma}\sum_{f=1}^{F}\sum_{j=1}^{F}\Gamma_{fj}^2.
\end{aligned}
\end{equation}
The first contribution measures disagreement with the mechanistic effluent state. The second measures violation of the modeled invariant relations. The third measures disagreement between the auxiliary fitted state and the coupled regression equation. The final two terms penalize the magnitudes of the driver and coupling coefficients.

The same structure can be written compactly in matrix form. The entry in row \(i\), column \(k\) of \((\widehat C-C_{in})A^\top\) is the invariant residual for observation \(i\). Likewise, the entry in row \(i\), column \(f\) of \(\widehat C R^\top\) is \(\widehat c_{if}-\sum_j\Gamma_{fj}\widehat c_{ij}\). The transpose appears because observations are stored as rows, whereas the single-observation relation \(Rc=r\) uses column vectors.

For \(N\) training observations, \(\Phi\in\mathbb{R}^{N\times D}\) stores the feature vectors by row. The matrices \(C_{in}\) and \(C_{out}\) contain the corresponding influent and mechanistic effluent states. The auxiliary fitted-state matrix is \(\widehat C\in\mathbb{R}_{+}^{N\times F}\). The resulting criterion is
\begin{equation}
\equationname{Matrix form of the structured-regression training objective}
\label{eq:icsor_training_objective}
\begin{aligned}
J(B,\Gamma,\widehat C)={}&
\|C_{out}-\widehat C\|_F^2
+\lambda_{inv}\| (\widehat C-C_{in})A^\top\|_F^2\\
&+\lambda_{sys}\|\widehat C(I_F-\Gamma)^\top-\Phi B^\top\|_F^2\\
&+\lambda_B\|B\|_F^2+\lambda_{\Gamma}\|\Gamma\|_F^2.
\end{aligned}
\end{equation}
The estimation problem minimizes this criterion subject to non-negative fitted training states, a zero coupling diagonal, coupling box bounds, and the conditioning safeguard. All penalty weights are non-negative, while the adopted ridge weights and coupled-system weight are strictly positive.

The first term favors agreement with the mechanistic component response. The second guides the fitted states toward the adopted mass conservation relations. The third connects those states to the simultaneous regression equation. The final two terms regularize the coefficients. The auxiliary state therefore mediates among statistical fit, physical guidance, and consistency with the learned coupled representation.

A one-observation example makes these contributions concrete. Take \(c_{out}=(3,1)^\top\), \(c_{in}=(2,2)^\top\), \(\widehat c=(2.5,1.2)^\top\), and \(A=[1\ \ 1]\). Use the coupling matrix in Equation~\eqref{eq:icsor_small_coupling_matrices}, and for a single baseline feature take \(B=(2,1)^\top\), giving the driver \((2,1)^\top\). The five unweighted contributions are
\begin{equation}
\equationname{Illustrative numerical contributions to the training penalties}
\label{eq:icsor_penalty_numbers}
\begin{aligned}
\text{prediction error}&=(3-2.5)^2+(1-1.2)^2=0.29,\\
\text{inventory error}&=(2.5+1.2-2-2)^2=0.09,\\
\text{system error}&=(2.5-0.2(1.2)-2)^2\\
&\quad +(1.2-0.2(2.5)-1)^2=0.1576,\\
\text{driver squared sum}&=2^2+1^2=5,\\
\text{coupling squared sum}&=0.2^2+0.2^2=0.08.
\end{aligned}
\end{equation}
With illustrative weights \(\lambda_{inv}=2\), \(\lambda_{sys}=3\), \(\lambda_B=0.1\), and \(\lambda_{\Gamma}=0.5\), the objective becomes \(0.29+2(0.09)+3(0.1576)+0.1(5)+0.5(0.08)=1.4828\). These illustrative weights simply show how the objective assigns relative importance to the different discrepancies. The numerical values used in the actual fit are reported in Table~\ref{tab:icsor_fitting_settings}.

The auxiliary state exists only within the fitting problem. At deployment, a new prediction is generated by solving the learned coupled relation. The invariant penalty therefore provides guidance during fitting rather than a final physical guarantee. This use of residual penalties follows the general physics-informed training principle developed by \citet{Raissi2019}. Mass conservation and non-negativity of the returned prediction are subsequently enforced by checked deployment.

The objective is formulated in physical component coordinates. Because the components have different units and numerical scales, the magnitudes of their residuals and fitted coefficients reflect those conventions. The penalty weights likewise depend on the adopted feature and invariant scaling. Any consistent rescaling must transform the features, coefficients, invariant operator, reporting map, and penalty weights together. \citet{Pinheiro2025} discuss the influence of feature scaling on learning. In this setting, physical interpretability requires both numerical conditioning and unit consistency.

\subsubsection{Deterministic block updates}

The product \(\widehat C\Gamma^\top\) makes the complete objective nonconvex. However, fixing two variable blocks at a time leaves a simpler problem in the remaining block. The fitting procedure therefore updates the driver coefficients, the coupling matrix, and the fitted states sequentially within each cycle.

The source of the nonconvexity can be seen from the reduced residual \((1-\gamma z)^2\), where \(\gamma\) is a coupling coefficient and \(z\) a fitted component. With \(z\) fixed, the expression is quadratic in \(\gamma\); with \(\gamma\) fixed, it is quadratic in \(z\). Jointly, however, it is not convex. At \((\gamma,z)=(0.2,5)\) and \((0.1,10)\), the residual is zero, while at the midpoint \((0.15,7.5)\) it becomes \((1-1.125)^2=1/64\). A convex function could not exceed the average of two zero endpoint values at their midpoint. The coupling between a fitted state and a learned coefficient therefore creates the nonconvex joint problem.

For fixed \(\Gamma\) and \(\widehat C\), the driver update is a ridge-regularized linear least-squares problem. Define \(T=\widehat C R^\top\). For one output row, let its coefficient vector be \(\beta\) and the corresponding target column of \(T\) be \(t\). The relevant part of the objective is
\begin{equation}
\equationname{Regularized least-squares objective for one driver coefficient row}
\label{eq:icsor_ridge_row_objective}
\lambda_{sys}\sum_{i=1}^{N}
\left(t_i-\sum_{\ell=1}^{D}\phi_{i\ell}\beta_\ell\right)^2
+\lambda_B\sum_{\ell=1}^{D}\beta_\ell^2.
\end{equation}
Differentiating with respect to coefficient \(\beta_h\) and setting the derivative to zero gives
\begin{equation}
\equationname{Scalar stationarity conditions for the driver coefficient row}
\label{eq:icsor_ridge_scalar_normal}
\lambda_{sys}\sum_{i=1}^{N}\phi_{ih}
\left(\sum_{\ell=1}^{D}\phi_{i\ell}\beta_\ell-t_i\right)
+\lambda_B\beta_h=0.
\end{equation}
Writing the condition for all coefficients produces the normal equations
\begin{equation}
\equationname{Normal equations for the driver coefficient row}
\label{eq:icsor_ridge_normal_equations}
\left(\lambda_{sys}\Phi^\top\Phi+\lambda_BI_D\right)\beta
=\lambda_{sys}\Phi^\top t.
\end{equation}
The matrix \(\Phi^\top\Phi\) contains the pairwise products of feature columns, while \(\Phi^\top t\) contains the products of features with the target. The positive ridge weight adds curvature in directions that would otherwise be weakly identified, including those created by duplicated ordered quadratic columns.

For example, with one coefficient and no intercept, take \(\Phi=(1,2)^\top\), \(t=(2,3)^\top\), and unit weights. The objective is
\begin{equation}
\equationname{Illustrative scalar ridge objective and its polynomial expansion}
\label{eq:icsor_ridge_numeric_example}
(2-\beta)^2+(3-2\beta)^2+\beta^2
=13-16\beta+6\beta^2.
\end{equation}
Its derivative is \(-16+12\beta\), so the minimizer is \(\beta=4/3\). The matrix form gives \(\Phi^\top\Phi=5\), \(\Phi^\top t=8\), and \((5+1)\beta=8\), yielding the same result. Without the ridge term, the coefficient would be \(8/5\). The example therefore illustrates how regularization shrinks the fitted coefficient.

For all output rows, the driver update becomes
\begin{equation}
\equationname{Closed-form ridge update for the driver coefficient matrix}
\label{eq:icsor_ridge_update}
B^\top=
\left(\lambda_{sys}\Phi^\top\Phi+\lambda_B I_D\right)^{-1}
\lambda_{sys}\Phi^\top\widehat C R^\top.
\end{equation}
The positive driver ridge weight makes the coefficient matrix positive definite even when some feature columns are duplicated by the ordered quadratic expansion. In computation, the corresponding linear system is solved directly rather than forming the explicit matrix inverse. The square quadratic slices are then symmetrized without changing the fitted response.

For fixed \(B\) and \(\widehat C\), define \(E=\widehat C-\Phi B^\top\). The coupling calculation can then be read one output row at a time. For row \(f\), the residual at observation \(i\) is \(E_{if}-\sum_{j\neq f}\widehat c_{ij}\Gamma_{fj}\). The other fitted component columns therefore act as predictors of \(E_{if}\), subject to ridge regularization and box constraints.

A reduced example illustrates the update. Suppose one permitted coupling \(\gamma\) uses predictor values 1 and 2 and residual targets 0.4 and 0.8. With equal unit weights,
\begin{equation}
\equationname{Illustrative scalar coupling objective and its expansion}
\label{eq:icsor_coupling_numeric_update}
(0.4-\gamma)^2+(0.8-2\gamma)^2+\gamma^2
=0.8-4\gamma+6\gamma^2.
\end{equation}
The derivative \(-4+12\gamma\) gives \(\gamma=1/3\), which lies below the adopted absolute entry bound of 0.368. If the box bound were instead 0.25, the constrained optimum would occur at that boundary. The full conditioning test is then applied to the assembled coupling matrix rather than to each coefficient separately.

Collecting the row-wise problems gives
\begin{equation}
\equationname{Regularized least-squares update for component coupling}
\label{eq:icsor_coupling_update}
\min_{\Gamma}\quad
\lambda_{sys}\|E-\widehat C\Gamma^\top\|_F^2+
\lambda_{\Gamma}\|\Gamma\|_F^2
\end{equation}
subject to the zero diagonal and box bounds. This subproblem is a convex quadratic program. Its candidate is subsequently screened using Equation~\eqref{eq:icsor_coupling_safeguards}. If the conditioning requirement is violated, the candidate is shrunk toward zero or replaced by an admissible fallback. The screened proposal is then compared with the previous admissible coupling matrix using the complete objective.

For fixed \(B\) and \(\Gamma\), the fitted-state update separates across observations. For sample \(i\), write \(\widehat c_i\) as a column vector. Expanding the three state-dependent residual terms gives a quadratic function of \(\widehat c_i\). The data term contributes \(\widehat c^\top\widehat c-2c_{out}^\top\widehat c\), the invariant term contributes \(\lambda_{inv}\widehat c^\top A^\top A\widehat c-2\lambda_{inv}c_{in}^\top A^\top A\widehat c\), and the system term contributes \(\lambda_{sys}\widehat c^\top R^\top R\widehat c-2\lambda_{sys}(B\phi)^\top R\widehat c\). Terms independent of \(\widehat c\) do not affect the minimizer.

Writing the resulting objective as \(\widehat c_i^\top H_C\widehat c_i-2h_i^\top\widehat c_i\), the matrix entries are
\begin{equation}
\equationname{Scalar entries of the fitted-state Hessian and forcing vector}
\label{eq:icsor_state_scalar_coefficients}
\begin{aligned}
(H_C)_{fg}&=\delta_{fg}+\lambda_{inv}\sum_{k=1}^{K}A_{kf}A_{kg}
+\lambda_{sys}\sum_{j=1}^{F}R_{jf}R_{jg},\\
(h_i)_f&=c_{out,if}
+\lambda_{inv}\sum_{k=1}^{K}A_{kf}\sum_{g=1}^{F}A_{kg}c_{in,ig}\\
&\quad+\lambda_{sys}\sum_{j=1}^{F}R_{jf}\sum_{\ell=1}^{D}B_{j\ell}\phi_{i\ell}.
\end{aligned}
\end{equation}
Here \(\delta_{fg}\) is one when \(f=g\) and zero otherwise. The off-diagonal terms show that the preferred value of one fitted component can depend on changes in another through both the invariant and coupled-system contributions.

Completing the square clarifies what the non-negative solve is doing. With \(c_0=H_C^{-1}h_i\),
\begin{equation}
\equationname{Completed-square form of the fitted-state objective}
\label{eq:icsor_state_completed_square}
\widehat c^\top H_C\widehat c-2h_i^\top\widehat c
=(\widehat c-c_0)^\top H_C(\widehat c-c_0)
-h_i^\top H_C^{-1}h_i.
\end{equation}
The unrestricted minimizer is \(c_0\). If one or more entries are negative, the constrained problem finds the closest non-negative state in the quadratic geometry defined by \(H_C\). Because that geometry can contain off-diagonal coupling, simple coordinate clipping need not yield the constrained optimum.

For example, take
\(H_C=\begin{bmatrix}3&1\\1&3\end{bmatrix}\) and \(h_i=(-1,2)^\top\). These coefficients can arise from unit data, invariant, and system weights with \(A=[1\ \ 1]\) and \(R=I\). The unrestricted solution is \(c_0=(-5/8,7/8)^\top\). Componentwise clipping gives \((0,7/8)^\top\). If the first component is instead fixed at zero and the remaining quadratic \(3c_2^2-4c_2\) is minimized, the optimum is \(c_2=2/3\).

At \((0,2/3)^\top\), the derivative in the first coordinate is \(2(c_2+1)=10/3\), so increasing the component away from its zero boundary would increase the objective. The objective value is \(-4/3\) at this constrained solution and \(-77/64\) at the clipped state, confirming that clipping is not optimal. These negative values arise because constant terms were omitted during the quadratic expansion; the complete residual-based training criterion remains non-negative.

The general fitted-state coefficients are
\begin{equation}
\equationname{Hessian and forcing vector for the fitted-state update}
\label{eq:icsor_state_update}
\begin{aligned}
H_C&=I_F+\lambda_{inv}A^\top A+\lambda_{sys}R^\top R,\\
h_i&=c_{out,i}+\lambda_{inv}A^\top A c_{in,i}
+\lambda_{sys}R^\top B\phi_i.
\end{aligned}
\end{equation}
Each observation minimizes \(\widehat c_i^\top H_C\widehat c_i-2h_i^\top\widehat c_i\) over \(\widehat c_i\geq\mathbf{0}\). The identity contribution makes \(H_C\) positive definite, so the row minimizer is unique. The off-diagonal terms again explain why componentwise clipping of the unconstrained solution generally does not solve the constrained problem.

The quadratic-program method of \citet{Stellato2018} provides an algorithmic basis for the constrained coupling and fitted-state updates. A generic stochastic first-order method could also optimize a related objective, but additional mechanisms would still be required to control non-negativity, coupling bounds, and matrix conditioning. The block formulation keeps each requirement associated with a clearly defined subproblem.

\subsubsection{Descent, stopping, and numerical safeguards}

Exact minimization of one block cannot increase the objective when the previous block value remains feasible. Since \(J\geq0\), a monotone non-increasing objective sequence converges in value. The convergence analysis of \citet{Razaviyayn2013} relates blockwise minimization to stationary solutions under its stated assumptions.

The conditioning screen requires an additional safeguard because shrinking a proposed coupling matrix can increase the objective relative to the previous admissible matrix. The procedure therefore evaluates the complete objective after screening and retains the previous admissible coupling whenever the screened candidate does not provide an acceptable non-increasing update. Numerical solver tolerances are also allowed explicitly. Monotone descent therefore applies to the accepted sequence of screened objective-decreasing updates.

Each restart begins from an admissible coupling matrix and non-negative fitted states. A ridge solve provides the initial driver. Complete cycles then continue until the iteration limit is reached or the slope of the recent objective history is sufficiently small. A running minimum records the best admissible state encountered. If several restarts are used, the admissible state with the smallest recorded objective is retained. A nearly flat objective therefore triggers stopping without discarding a better state reached earlier in the sequence.

Figure~\ref{fig:ch5_training_cycle} separates the three fitting updates from the checks used to accept them. The reduced calculations in Equations~\eqref{eq:icsor_ridge_numeric_example} and \eqref{eq:icsor_coupling_numeric_update} illustrate the driver and coupling updates, while the non-negative-state example explains why the fitted-state block requires a constrained solve.

\begin{figure}[htbp]
\centering
\begingroup\linespread{1}\selectfont
\begin{tikzpicture}[
    c5box/.style={draw=blue!50!black,fill=blue!5,rounded corners=2pt,
        line width=0.7pt,align=center,font=\small,inner sep=6pt},
    c5decision/.style={diamond,aspect=2.2,draw=blue!50!black,fill=black!4,
        line width=0.7pt,align=center,font=\small,text width=2.1cm,inner sep=2pt},
    c5arrow/.style={-{Latex[length=2.2mm]},draw=blue!65!black,line width=0.8pt},
    c5label/.style={font=\small,fill=white,inner sep=2pt}
]
\node[c5box,text width=4.5cm] (initialize) at (0,0)
    {\textbf{Start a restart}\\Admissible \(\Gamma\) and \(\widehat C\geq0\)};
\node[c5box,text width=4.5cm,below=6mm of initialize] (driver)
    {\textbf{Driver ridge update}\\Fit \(B\) with \(\Gamma,\widehat C\) fixed};
\node[c5box,text width=4.5cm,below=6mm of driver] (symmetry)
    {\textbf{Symmetric quadratic slices}\\Keep predictions and reduce the ridge norm};
\node[c5box,text width=4.5cm,below=6mm of symmetry] (proposal)
    {\textbf{Coupling proposal}\\Zero diagonal and box bounds};
\node[c5box,text width=4.0cm,right=12mm of proposal] (screen)
    {\textbf{Screen the proposal}\\Check conditioning\\and objective\\Shrink or retain \(\Gamma\)};
\node[c5box,text width=4.0cm,above=6mm of screen] (states)
    {\textbf{Non-negative fitted states}\\Update \(\widehat C\) with \(B,\Gamma\) fixed};
\node[c5box,text width=4.0cm,above=6mm of states] (record)
    {\textbf{Evaluate the objective}\\Record the best\\admissible objective};
\node[c5decision,above=6mm of record] (stop)
    {Small slope\\or limit?};
\node[c5box,text width=4.0cm,above=6mm of stop] (retain)
    {\textbf{Complete the restart}\\Retain its best state};
\node[c5decision] (restart) at (initialize.center |- retain.center)
    {More planned\\restarts?};
\node[c5box,text width=2.25cm,left=7mm of restart] (finish)
    {\textbf{Best fit}\\Smallest admissible objective};
\coordinate (loopcorner) at ($(initialize.east)+(5mm,0)$);
\draw[c5arrow] (initialize.south)--(driver.north);
\draw[c5arrow] (driver.south)--(symmetry.north);
\draw[c5arrow] (symmetry.south)--(proposal.north);
\draw[c5arrow] (proposal.east)--(screen.west);
\draw[c5arrow] (screen.north)--(states.south);
\draw[c5arrow] (states.north)--(record.south);
\draw[c5arrow] (record.north)--(stop.south);
\draw[c5arrow] (stop.north)--node[c5label,right] {Yes}(retain.south);
\draw[c5arrow] (stop.west)--node[c5label,above] {No}(loopcorner |- stop.west)
    --(loopcorner |- driver.east)--(driver.east);
\draw[c5arrow] (retain.west)--(restart.east);
\draw[c5arrow] (restart.south)--node[c5label,right] {Yes}(initialize.north);
\draw[c5arrow] (restart.west)--node[c5label,above] {No}(finish.east);
\end{tikzpicture}
\endgroup

\caption{Conceptual fitting cycle with explicit symmetry, conditioning, descent, and non-negativity checks. The prescribed restart count can be one.}
\label{fig:ch5_training_cycle}
\end{figure}

The cycle returns to the driver because each block changes the conditions faced by the remaining blocks. The conditioning screen is applied to the coupling proposal before the fitted-state update, and the complete objective is recorded after each full cycle. The restart decision belongs to the prescribed search procedure, while the final model selection compares the admissible objective values obtained. Together, these steps define the accepted sequence of updates for the nonconvex fit.

\begin{table}[htbp]
\centering
\caption{Specified fitting settings for the structured regression.}
\label{tab:icsor_fitting_settings}
\begin{tabular}{ll}
\toprule
Setting & Value\\
\midrule
Baseline feature & Included\\
Invariant penalty \(\lambda_{inv}\) & \(8.68\times10^{-3}\)\\
Coupled-system penalty \(\lambda_{sys}\) & \(1.27\times10^{-3}\)\\
Driver ridge penalty \(\lambda_B\) & \(6.09\times10^{-5}\)\\
Coupling ridge penalty \(\lambda_{\Gamma}\) & \(6.55\times10^{-4}\)\\
Off-diagonal bound \(\gamma_{abs}\) & 0.368\\
Maximum full fitting cycles & 60\\
Restarts & 1\\
Trailing objective window & 12 cycles\\
Objective slope tolerance & \(2.85\times10^{-7}\)\\
\bottomrule
\end{tabular}
\end{table}

Table~\ref{tab:icsor_fitting_settings} summarizes the numerical choices used during fitting. The condition threshold \(\kappa_{max}\) is an additional prescribed safeguard and is checked independently of the coefficient bound. Hyperparameter assessment uses validation observations drawn only from the training pool. Held-out test responses do not determine the penalty weights or stopping settings \citep{Kaufman2012}. The fitted settings are held fixed within each sample-size comparison so that changes in the learning curves reflect differences in available data rather than changes in the model specification.

\subsection{Checked deployment}

\subsubsection{The raw coupled state}

For a new input pair \((u_*,c_{in,*})\), the fitted parameters define
\begin{equation}
\equationname{Raw deployed state from the fitted coupled system}
\label{eq:icsor_raw_state}
\widehat R=I_F-\widehat\Gamma,
\qquad d_* =\widehat B\phi(u_*,c_{in,*}),
\qquad \widehat R c_{raw}=d_*.
\end{equation}
The notation \(c_{raw}=\widehat R^{-1}d_*\) describes the mapping mathematically, although the prediction is evaluated numerically by solving the linear system. The adopted conditioning bound limits the sensitivity of that solve. The raw state can be returned directly only if it satisfies both the invariant and non-negativity checks.

The two diagnostics measure different departures from physical admissibility. For a vector, the Euclidean norm is the square root of the sum of squared entries, so an invariant-residual vector \((3,4)^\top\) has norm 5. For an illustrative concentration vector \((-0.3,0.1,-0.4)^\top\), retaining only the negative entries gives \((-0.3,0,-0.4)^\top\), whose absolute norm is 0.7. The first quantity measures the size of a mass conservation discrepancy; the second measures the total negative concentration in the adopted numerical coordinates.

For a candidate state \(c\), define
\begin{equation}
\equationname{Mass conservation and non-negativity diagnostics of a deployed state}
\label{eq:icsor_deployment_diagnostics}
v_{mc}(c)=\|A(c-c_{in,*})\|_2,
\qquad
v_{nn}(c)=\|\min(c,\mathbf{0})\|_1.
\end{equation}
The minimum is evaluated componentwise. The numerical diagnostic tolerance is \(10^{-10}\). The mathematical requirements are exact equality and non-negativity; the residual tolerance defines their computational assessment under the stated scaling of \(A\).

\subsubsection{Affine projection for mass conservation}

The first correction stage enforces the mass conservation equalities without yet imposing concentration bounds. The underlying geometry can again be seen in the two-component example. Let the raw prediction be \((p_1,p_2)^\top\) and let the required total inventory be \(T\). The nearest mass-conserving state minimizes \((c_1-p_1)^2+(c_2-p_2)^2\) subject to \(c_1+c_2=T\). Introducing multiplier \(\eta\) gives
\begin{equation}
\equationname{Lagrangian for an illustrative two-component affine projection}
\label{eq:icsor_scalar_projection_lagrangian}
\mathcal{L}=\frac{1}{2}(c_1-p_1)^2+
\frac{1}{2}(c_2-p_2)^2+\eta(c_1+c_2-T).
\end{equation}
The stationarity and equality conditions are
\begin{equation}
\equationname{Stationarity and equality conditions for the illustrative affine projection}
\label{eq:icsor_scalar_projection_conditions}
c_1-p_1+\eta=0,
\qquad c_2-p_2+\eta=0,
\qquad c_1+c_2=T.
\end{equation}
Thus
\begin{equation}
\equationname{Closed-form solution of the illustrative affine projection}
\label{eq:icsor_scalar_projection_solution}
\begin{aligned}
\eta&=\frac{p_1+p_2-T}{2},\\
c_1&=p_1-\frac{p_1+p_2-T}{2},\\
c_2&=p_2-\frac{p_1+p_2-T}{2}.
\end{aligned}
\end{equation}
The unweighted projection distributes the inventory discrepancy equally between the two numerical coordinates. For \(p_1=12\), \(p_2=1\), and \(T=10\), the discrepancy is 3, so the multiplier is 1.5 and the projected state is \((10.5,-0.5)^\top\). The total inventory is correct, but the second concentration is negative. This example shows why an equality-only projection cannot guarantee complete physical admissibility.

In the full component space, a raw state that fails the checks is first projected by
\begin{equation}
\equationname{Euclidean affine projection onto the mass conservation equality}
\label{eq:icsor_affine_problem}
c_{aff}=\mathop{\arg\min}_{c\in\mathbb{R}^{F}}\|c-c_{raw}\|_2^2
\quad\text{subject to}\quad Ac=Ac_{in,*}.
\end{equation}
With multiplier \(\eta\in\mathbb{R}^{K}\), stationarity gives \(c-c_{raw}+A^\top\eta=0\). Substitution into the equality yields \(AA^\top\eta=A(c_{raw}-c_{in,*})\), and full row rank of \(A\) gives
\begin{equation}
\equationname{Closed-form affine projection of the raw component state}
\label{eq:icsor_affine_solution}
c_{aff}=c_{raw}-A^\top(AA^\top)^{-1}A(c_{raw}-c_{in,*}).
\end{equation}
Because \(A\) is fixed, the matrix \(A^\top(AA^\top)^{-1}A\) can be prepared once and reused for every new prediction.

The two-component example makes the same matrix operation explicit:
\begin{equation}
\equationname{Illustrative matrix of the mass conservation projection}
\label{eq:icsor_small_affine_matrix}
AA^\top=2,
\qquad
A^\top(AA^\top)^{-1}A
=\frac{1}{2}\begin{bmatrix}1&1\\1&1\end{bmatrix}.
\end{equation}
With illustrative influent \((6,4)^\top\), the raw-minus-influent difference is \((6,-3)^\top\). Multiplication by the projection matrix gives \((1.5,1.5)^\top\), which is subtracted from the raw state. The matrix calculation therefore reproduces the scalar result without requiring an arbitrary choice about which component should absorb the imbalance.

Geometrically, this is the orthogonal projection onto the affine set defined by the adopted mass conservation equations in the chosen numerical component coordinates. It removes the part of the raw prediction normal to that set while retaining the tangent component. Because the distance is measured in the stated component units, the scaling of those coordinates affects the adjustment. \citet{SturmSilva2024} develop related mass-conservation projections and show why species-dependent weighting matters when concentration scales differ.

The affine state is checked for both balance residuals and non-negativity before it is accepted. The equality follows analytically from Equation~\eqref{eq:icsor_affine_solution}, but the numerical residual check can still reveal inconsistency in the operator or computation. If the affine state remains negative, a second projection must enforce the equality and inequality requirements simultaneously.

\subsubsection{Joint mass conservation and non-negativity}

The final projection can again be understood from the two-component example. Starting from the mass-conserving but negative state \((10.5,-0.5)^\top\), require \(c_1+c_2=10\) and \(c_1,c_2\geq0\). Introduce non-negative adjustment variables \(\rho_1,\rho_2\) satisfying
\begin{equation}
\equationname{Absolute-adjustment bounds for an illustrative two-component projection}
\label{eq:icsor_scalar_absolute_constraints}
\begin{aligned}
\rho_1&\geq c_1-10.5,&\rho_1&\geq10.5-c_1,\\
\rho_2&\geq c_2+0.5,&\rho_2&\geq-0.5-c_2.
\end{aligned}
\end{equation}
These inequalities make \(\rho_1\) and \(\rho_2\) at least as large as the absolute adjustments of the two coordinates. Minimizing \(w_1\rho_1+w_2\rho_2\) with positive weights makes them equal to those absolute adjustments at the optimum.

Imposing the inventory equality with \(c_2=t\) gives \(c_1=10-t\) and \(0\leq t\leq10\). Both absolute adjustments are then \(t+0.5\), so the objective becomes \((w_1+w_2)(t+0.5)\). It is minimized at \(t=0\), yielding \((10,0)^\top\). With equal weights, this state has cost 1, whereas \(t=1\) would give \((9,1)^\top\) with cost 3. The reduced example therefore makes the optimal joint correction transparent.

For the full component state, the final projection is
\begin{equation}
\equationname{Weighted absolute-distance projection for mass conservation and non-negativity}
\label{eq:icsor_weighted_correction}
\begin{aligned}
\widehat c =\mathop{\arg\min}_{c,\rho\in\mathbb{R}^{F}}\quad
& w_c^\top\rho\\
\text{subject to}\quad
& Ac=Ac_{in,*},\\
& c\geq\mathbf{0},\\
& -\rho\leq c-c_{aff}\leq\rho,\\
& \rho\geq\mathbf{0}.
\end{aligned}
\end{equation}
For strictly positive weights \(w_c\), the optimal auxiliary variable satisfies \(\rho_f=|c_f-c_{aff,f}|\). The objective is therefore the weighted absolute adjustment \(\sum_f w_{c,f}|c_f-c_{aff,f}|\). The fitted regression itself remains unchanged; only the returned state is corrected.

The problem is linear because the absolute values are represented by auxiliary variables and linear inequalities. The revised-simplex method described by \citet{HuangfuHall2018} provides an algorithmic basis for this type of projection. The mathematical feasible set is nonempty whenever the supplied influent state is non-negative and satisfies the adopted boundary relation. Positive weights make the objective bounded below and discourage unnecessary movement in every component. A minimizer therefore exists, although several minimizing states can occur if the weighted absolute-distance objective contains a flat face.

The weights determine how the required adjustment is distributed. A large weight makes a component expensive to change, whereas a smaller weight allows more movement if that helps satisfy the constraints. The weights are fixed before deployment according to the intended component scaling and scientific priorities. Equal numerical weights are themselves a modeling choice because the components differ in units and typical magnitudes.

A three-component illustration shows how this choice matters. Let the invariant row be \([1\ \ 1\ \ 1]\), and let the affine reference be \((-1,4,7)^\top\), whose total is 10. Restoring the first component to zero requires one unit to be removed elsewhere. Two feasible candidates are \((0,3,7)^\top\) and \((0,4,6)^\top\). With weights \((1,10,1)^\top\), their costs are 11 and 2, respectively, so the second candidate is preferred because changing the second component is expensive. With weights \((1,1,10)^\top\), the costs reverse to 2 and 11. Both candidates conserve the same total inventory; the weights determine which admissible redistribution is preferred.

\citet{KircherVotsmeier2025} show why mass-conservation projection in reaction-kinetics surrogates must be combined with non-negativity enforcement. Their scaled-space formulation protects low-concentration species during correction. Equation~\eqref{eq:icsor_weighted_correction} addresses the same joint requirement for activated sludge components through a weighted absolute-distance formulation. It enforces the adopted linear invariants and non-negativity on the final returned state.

\subsubsection{Illustrative projection and the failure of clipping}

The difference between projection and simple clipping can be summarized using the same two-component system. Let \(A=[1\ \ 1]\), let the influent inventory total be 10, and let the raw prediction be \(c_{raw}=(12,1)^\top\). The raw total is 13. The affine projection subtracts 1.5 from each coordinate, producing \(c_{aff}=(10.5,-0.5)^\top\). Mass conservation is restored, but the second component becomes negative.

If that negative concentration is clipped to zero, the state becomes \((10.5,0)^\top\), whose total is 10.5 and therefore violates the invariant. Equation~\eqref{eq:icsor_weighted_correction}, by contrast, returns \((10,0)^\top\) for equal weights. Both requirements are satisfied simultaneously. The distinction is important because repairing one physical defect independently can recreate another.

The deployment procedure therefore follows the geometry of the feasible set. A raw state that already satisfies both checks is returned without modification. Otherwise, the affine projection first restores the invariant relations. If that state is also non-negative, it can be returned. The linear program is needed only when the affine state remains outside the non-negative region. Every returned candidate is then checked under the declared numerical tolerances, and a failed solve or failed final audit is recorded as a prediction failure.

Figure~\ref{fig:ch5_checked_deployment} makes these branches explicit.

\begin{figure}[htbp]
\centering
\begingroup\linespread{1}\selectfont
\begin{tikzpicture}[
    c5box/.style={draw=blue!50!black,fill=blue!5,rounded corners=2pt,
        line width=0.7pt,align=center,font=\small,inner sep=6pt},
    c5decision/.style={diamond,aspect=2.3,draw=blue!50!black,fill=black!4,
        line width=0.7pt,align=center,font=\small,text width=2.3cm,inner sep=2pt},
    c5arrow/.style={-{Latex[length=2.2mm]},draw=blue!65!black,line width=0.8pt},
    c5label/.style={font=\small,fill=white,inner sep=2pt}
]
\node[c5box,text width=4.8cm] (raw) at (0,0)
    {\textbf{Raw coupled prediction}\\\(c_{raw}=\widehat R^{-1}d_*\)};
\node[c5decision,below=6mm of raw] (rawcheck)
    {Both checks\\pass?};
\node[c5box,text width=3.3cm] (rawaccept) at (6,0 |- rawcheck.center)
    {Set \(\widehat c=c_{raw}\)};
\node[c5box,text width=4.8cm,below=6mm of rawcheck] (affine)
    {\textbf{Project for mass conservation}\\Compute \(c_{aff}\)};
\node[c5decision,below=6mm of affine] (affcheck)
    {Both checks\\pass?};
\node[c5box,text width=3.3cm] (affaccept) at (6,0 |- affcheck.center)
    {Set \(\widehat c=c_{aff}\)};
\node[c5box,text width=4.8cm,below=6mm of affcheck] (projection)
    {\textbf{Weighted projection}\\Minimize \(w_c^\top\rho\)\\\(Ac=Ac_{in,*},\quad c\geq0\)\\\(\rho\geq|c-c_{aff}|\)};
\node[c5decision,below=6mm of projection] (audit)
    {Solve and\\audit pass?};
\node[c5box,text width=3.3cm] (lpaccept) at (6,0 |- audit.center)
    {Use optimized \(\widehat c\)};
\node[c5box,text width=4.8cm,fill=black!5,draw=black!60,below=6mm of audit] (failed)
    {\textbf{Prediction failure}\\Admissibility not verified};
\node[c5box,text width=3.7cm] (return) at (6,0 |- failed.center)
    {\textbf{Checked output}\\\(\widehat c,\quad I_{\mathrm{comp}}\widehat c\)};
\coordinate (merge) at (8.5,0);
\draw[c5arrow] (raw.south)--(rawcheck.north);
\draw[c5arrow] (rawcheck.east)--node[c5label,above] {Yes}(rawaccept.west);
\draw[c5arrow] (rawcheck.south)--node[c5label,right] {No}(affine.north);
\draw[c5arrow] (affine.south)--(affcheck.north);
\draw[c5arrow] (affcheck.east)--node[c5label,above] {Yes}(affaccept.west);
\draw[c5arrow] (affcheck.south)--node[c5label,right] {No}(projection.north);
\draw[c5arrow] (projection.south)--(audit.north);
\draw[c5arrow] (audit.east)--node[c5label,above] {Yes}(lpaccept.west);
\draw[c5arrow] (audit.south)--node[c5label,right] {No}(failed.north);
\draw[c5arrow] (rawaccept.east)--(merge |- rawaccept.east)
    --(merge |- return.east)--(return.east);
\draw[draw=blue!65!black,line width=0.8pt] (affaccept.east)--(merge |- affaccept.east);
\draw[draw=blue!65!black,line width=0.8pt] (lpaccept.east)--(merge |- lpaccept.east);
\fill[blue!65!black] (merge |- affaccept.east) circle (1.2pt);
\fill[blue!65!black] (merge |- lpaccept.east) circle (1.2pt);
\end{tikzpicture}
\endgroup

\caption{Checked deployment of one predicted component state. Accepted branches merge before the common reporting map is applied. A failed final audit produces a reported prediction failure.}
\label{fig:ch5_checked_deployment}
\end{figure}

The illustrative state \((12,1)^\top\) follows both ``No'' branches because the raw inventory is incorrect and the affine projection contains a negative concentration. The weighted projection produces \((10,0)^\top\), which enters the return branch only after both checks pass. A state that satisfies the requirements earlier in the sequence reaches the same reporting step with less projection effort. The branches therefore differ in computational work, but not in the physical acceptance requirements.

The final component state is the deployed model output. The corresponding reported water-quality quantities are
\begin{equation}
\equationname{Water-quality composites of the final checked state}
\label{eq:icsor_deployed_composites}
\widehat y_{\mathrm{comp}}=I_{\mathrm{comp}}\widehat c.
\end{equation}
The deployment constraints determine whether the state is admissible, while the stage-wise assessment records how often each correction stage is needed.

\subsection{Interpreting the complete prediction map}

\subsubsection{Driver coefficients and back projection}

The effect of component coupling on coefficient interpretation can be seen using the same two-component example. Take the coupling matrix from Equation~\eqref{eq:icsor_small_coupling_matrices}, a reduced feature vector \((1,a)^\top\), and driver coefficients
\begin{equation}
\equationname{Illustrative driver coefficients and operating-dependent driver}
\label{eq:icsor_small_driver_coefficient_matrix}
B=\begin{bmatrix}4&1\\2&0\end{bmatrix},
\qquad
r=\begin{bmatrix}4+a\\2\end{bmatrix}.
\end{equation}
The first coefficient column gives the baseline response, and the second gives direct dependence on the hypothetical scalar input \(a\). Passing these coefficients through the coupling inverse gives
\begin{equation}
\equationname{Illustrative transformation of driver coefficients through coupling}
\label{eq:icsor_small_back_projection}
\begin{aligned}
R^{-1}B
&=\frac{1}{0.96}
\begin{bmatrix}1&0.2\\0.2&1\end{bmatrix}
\begin{bmatrix}4&1\\2&0\end{bmatrix}\\
&=\frac{1}{0.96}\begin{bmatrix}4.4&1\\2.8&0.2\end{bmatrix}
=\begin{bmatrix}55/12&25/24\\35/12&5/24\end{bmatrix}.
\end{aligned}
\end{equation}
The second raw component now depends on \(a\), even though its own driver coefficient for that input was zero. The dependence is introduced by the simultaneous coupling solve. A coefficient in the driver is therefore not, by itself, the coefficient of the final raw component response.

The same transformation continues into reported water-quality quantities. For a hypothetical quantity \(y=c_1+2c_2\), with reporting row \([1\ \ 2]\),
\begin{equation}
\equationname{Illustrative composite coefficients after coupling}
\label{eq:icsor_small_composite_projection}
\begin{bmatrix}1&2\end{bmatrix}R^{-1}B
=\begin{bmatrix}125/12&35/24\end{bmatrix}.
\end{equation}
At \(a=2\), the driver is \((6,2)^\top\), the raw component state is \((20/3,10/3)^\top\), and the reported quantity is \(40/3\). The transformed coefficient row gives the same value through \(125/12+(35/24)2=40/3\). The example shows how the driver coefficients acquire their final raw-response meaning only after coupling and reporting have been applied.

For the fitted model, the raw component coefficient matrix is
\begin{equation}
\equationname{Raw component-response coefficients after coupling}
\label{eq:icsor_raw_component_coefficients}
B_{\mathrm{raw}}=\widehat R^{-1}\widehat B,
\qquad c_{raw}=B_{\mathrm{raw}}\phi.
\end{equation}
The corresponding coefficient map for COD, TN, TP, and TSS is
\begin{equation}
\equationname{Raw composite-response coefficients after coupling and reporting}
\label{eq:icsor_raw_composite_coefficients}
B_{\mathrm{raw,comp}}=I_{\mathrm{comp}}\widehat R^{-1}\widehat B,
\qquad y_{\mathrm{raw,comp}}=B_{\mathrm{raw,comp}}\phi.
\end{equation}
This transformation carries the driver coefficients first through the learned component coupling and then through the reporting map. Consequently, inspection of a COD quadratic term uses the COD row of \(B_{\mathrm{raw,comp}}\), rather than an arbitrary component row of \(\widehat B\).

Figure~\ref{fig:ch5_prediction_chain} summarizes this raw prediction chain while retaining the named feature blocks.

\begin{figure}[htbp]
\centering
\begingroup\linespread{1}\selectfont
\begin{tikzpicture}[
    c5box/.style={draw=blue!50!black,fill=blue!5,rounded corners=2pt,
        line width=0.7pt,align=center,font=\small,inner sep=7pt},
    c5feature/.style={c5box,text width=4.8cm,align=left},
    c5arrow/.style={-{Latex[length=2.2mm]},draw=blue!65!black,line width=0.8pt},
    c5label/.style={font=\small,fill=white,inner sep=2pt}
]
\node[c5box,text width=4.8cm] (inputs) at (0,0)
    {Operating conditions \(u\)\\Influent components \(c_{in}\)};
\node[c5feature,below=7mm of inputs] (features)
    {\textbf{Feature vector \(\phi\)}\\[3pt]
     Baseline \(1\)\\
     Operating \(u\)\\
     Influent \(c_{in}\)\\
     Operating pairs \(u\otimes u\)\\
     Influent pairs \(c_{in}\otimes c_{in}\)\\
     Operating-load pairs \(u\otimes c_{in}\)};
\node[c5box,text width=3.3cm] (driver) at ($(features.north east)+(2.5cm,-0.65cm)$)
    {\textbf{Component driver}\\\(r=B\phi\)};
\node[c5box,text width=3.3cm,below=7mm of driver] (solve)
    {\textbf{Coupled solve}\\\(Rc_{raw}=r\)};
\node[c5box,text width=3.3cm,fill=black!4,below=7mm of solve] (coupling)
    {Learned coupling\\\(R=I_F-\Gamma\)};
\node[c5box,text width=3.2cm,right=15mm of solve] (components)
    {\textbf{Raw component}\\\textbf{state}\\\(c_{raw}\)};
\node[c5box,text width=3.2cm,above=10mm of components] (reported)
    {\textbf{Reporting map}\\\(I_{\mathrm{comp}}c_{raw}\)};
\draw[c5arrow] (inputs.south)--(features.north);
\draw[c5arrow] (features.east)--node[c5label,above,sloped] {\(\phi\)}(driver.west);
\draw[c5arrow] (driver.south)--(solve.north);
\draw[c5arrow] (coupling.north)--(solve.south);
\draw[c5arrow] (solve.east)--node[c5label,above] {\(c_{raw}\)}(components.west);
\draw[c5arrow] (components.north)--(reported.south);
\end{tikzpicture}
\endgroup

\caption{Named feature blocks and the transformation from driver to raw component and reported-composite predictions. Composing the three operations gives the raw reported coefficient map \(I_{\mathrm{comp}}R^{-1}B\).}
\label{fig:ch5_prediction_chain}
\end{figure}

In the two-feature illustration, the second driver has no direct coefficient for the scalar input, but the coupling step gives the corresponding component a nonzero response. The reporting map then combines both component responses. Reading the calculation in this order explains the back projection in Equations~\eqref{eq:icsor_small_back_projection} and \eqref{eq:icsor_small_composite_projection}. Checked deployment occurs afterward and may further modify the component state before the final reported quantities are calculated.

The fitted coefficients are displayed according to the same logical structure. Operating terms are separated from influent terms, symmetric quadratic maps display the mirrored cells consistently, and operating-load maps show which operating variables interact with which influent components. The coupling matrix is displayed separately because its meaning differs from that of the driver coefficients.

A post-fitting display threshold is used only to simplify presentation. For each reported-composite raw-regression map, the threshold is \(10^{-4}\) times the largest absolute coefficient. For the coupling matrix, it is \(10^{-4}\) times the largest absolute off-diagonal entry. The fitted model itself remains unchanged; the threshold controls only which entries are shown. Symmetric quadratic summaries count unique upper-triangular terms. A small coefficient can still contribute substantially when multiplied by a large feature, so interpretation considers coefficient magnitude together with the scale of the corresponding input and its contribution over the declared domain \citep{Selerio2026,Nazlabadi2021}.

\subsubsection{Analytical jacobians of the raw regression}

A fitted coefficient and a local response derivative are not the same quantity. The scalar driver in Equation~\eqref{eq:icsor_scalar_driver_example} makes this distinction explicit:
\begin{equation}
\equationname{Gradients of the illustrative scalar regression driver}
\label{eq:icsor_scalar_driver_derivatives}
\begin{aligned}
\frac{\partial r_1}{\partial u_1}&=2+0.2u_1+0.2u_2,\\
\frac{\partial r_1}{\partial u_2}&=-1+0.2u_1+0.25c_1,\\
\frac{\partial r_1}{\partial c_1}&=0.5+0.25u_2,
&\frac{\partial r_1}{\partial c_2}&=-0.1c_2.
\end{aligned}
\end{equation}
At \(u=(2,3)^\top\) and \(c=(4,5)^\top\), these derivatives are 3, 0.4, 1.25, and \(-0.5\). The derivative with respect to \(u_1\) is therefore 3 even though the linear coefficient of \(u_1\) is only 2. The difference comes from the operating curvature and interaction terms.

The second illustrative driver \(r_2=-u_1+u_2+0.2c_2\) has derivative \(-1\) with respect to \(u_1\). Differentiating the coupled equations while holding the learned coupling fixed gives
\[
\frac{\partial c_1}{\partial u_1}-0.2\frac{\partial c_2}{\partial u_1}=3,
\qquad
-0.2\frac{\partial c_1}{\partial u_1}+\frac{\partial c_2}{\partial u_1}=-1.
\]
The same inverse used for prediction therefore propagates the derivative:
\begin{equation}
\equationname{Illustrative propagation of an operating derivative through coupling}
\label{eq:icsor_small_chain_derivative}
\begin{bmatrix}\partial c_1/\partial u_1\\\partial c_2/\partial u_1\end{bmatrix}
=\frac{1}{0.96}\begin{bmatrix}1&0.2\\0.2&1\end{bmatrix}
\begin{bmatrix}3\\-1\end{bmatrix}
=\begin{bmatrix}35/12\\-5/12\end{bmatrix}.
\end{equation}
For the hypothetical reporting row \([1\ \ 2]\), the resulting reported derivative is \(35/12+2(-5/12)=25/12\). Differentiating the driver, passing the derivative through coupling, and applying the reporting weights are therefore distinct stages of the same sensitivity calculation.

Collecting several derivatives produces a Jacobian. At the illustrative operating point,
\begin{equation}
\equationname{Illustrative operating Jacobians before and after coupling}
\label{eq:icsor_small_operating_jacobian}
J_{r,u}=\begin{bmatrix}3&0.4\\-1&1\end{bmatrix},
\qquad
R^{-1}J_{r,u}=\begin{bmatrix}35/12&5/8\\-5/12&9/8\end{bmatrix}.
\end{equation}
The first column reproduces the derivative calculation above. The second is obtained by solving the same coupled system for derivative column \((0.4,1)^\top\). The Jacobian therefore collects the local response of the complete raw mapping to several inputs at once.

For component \(f\), reshape the operating-load coefficients as \(H_{uc}^{(f)}\in\mathbb{R}^{M_{op}\times F}\), so that the interaction term is \(u^\top H_{uc}^{(f)}c_{in}\). The component driver can then be written as
\begin{equation}
\equationname{Component driver expressed through symmetric quadratic matrices}
\label{eq:icsor_driver_component}
\begin{aligned}
r_f={}&b_f+(W_u)_{f,:}u+(W_{in})_{f,:}c_{in}
+u^\top Q_{uu}^{(f)}u\\
&+c_{in}^\top Q_{cc}^{(f)}c_{in}
+u^\top H_{uc}^{(f)}c_{in}.
\end{aligned}
\end{equation}
With symmetric quadratic slices, its gradients are
\begin{equation}
\equationname{Operating and influent gradients of a component driver}
\label{eq:icsor_driver_gradients}
\begin{aligned}
\nabla_u r_f&=(W_u)_{f,:}^\top+2Q_{uu}^{(f)}u
+H_{uc}^{(f)}c_{in},\\
\nabla_{c_{in}}r_f&=(W_{in})_{f,:}^\top+2Q_{cc}^{(f)}c_{in}
+(H_{uc}^{(f)})^\top u.
\end{aligned}
\end{equation}
These expressions follow directly from differentiation. For \(x^\top Qx\), the derivative collects both mirrored terms and reduces to \(2Qx\) when \(Q\) is symmetric. For \(u^\top Hc\), differentiation gives \(Hc\) with respect to \(u\) and \(H^\top u\) with respect to \(c\).

Stacking the component gradients gives the driver Jacobians \(J_{r,u}\in\mathbb{R}^{F\times M_{op}}\) and \(J_{r,c}\in\mathbb{R}^{F\times F}\). The corresponding raw component and composite Jacobians are
\begin{equation}
\equationname{Raw component and composite Jacobians after coupling}
\label{eq:icsor_raw_jacobians}
\begin{aligned}
J_{raw,u}&=\widehat R^{-1}J_{r,u},
&J_{raw,c}&=\widehat R^{-1}J_{r,c},\\
J_{raw,comp,u}&=I_{\mathrm{comp}}J_{raw,u},
&J_{raw,comp,c}&=I_{\mathrm{comp}}J_{raw,c}.
\end{aligned}
\end{equation}
These derivatives describe the complete fitted raw regression at a specified operating point, including the effects of curvature, interactions, coupling, and reporting. Projection-stage derivatives are considered separately below.

For example, an aeration sensitivity can be evaluated at low and high ammonium loading while other inputs are held fixed. A difference between the two derivatives indicates a loading-dependent association in the fitted response. Oxygen-transfer limitations and heterotrophic competition provide process context for interpreting such a pattern \citep{StenstromSong1991}. The process model therefore provides a plausibility reference for the learned relationship without converting that statistical relationship into a causal claim.

\subsubsection{Sensitivity after the affine projection}

Projection can substantially alter the local response of the fitted model. Return to the two-feature example in Equation~\eqref{eq:icsor_small_back_projection}. Its raw derivative with respect to \(a\) is \((25/24,5/24)^\top\), whose entries sum to \(5/4\). This raw response changes the total inventory. If the influent inventory is fixed, the affine projection must remove that total change.

For the two-component invariant,
\begin{equation}
\equationname{Illustrative operating derivative after affine projection}
\label{eq:icsor_small_projection_derivative}
\begin{aligned}
P_A&=\frac{1}{2}\begin{bmatrix}1&-1\\-1&1\end{bmatrix},\\
\frac{\partial c_{aff}}{\partial a}
&=P_A\begin{bmatrix}25/24\\5/24\end{bmatrix}
=\begin{bmatrix}5/12\\-5/12\end{bmatrix}.
\end{aligned}
\end{equation}
The projected derivatives now sum to zero, as required by the fixed total inventory. For reporting row \([1\ \ 2]\), the projected derivative is \(5/12+2(-5/12)=-5/12\), whereas the raw derivative was \(35/24\). Projection therefore changes both the magnitude and sign of the local reported response because it enforces a different material-accounting condition on the complete state.

Influent sensitivities involve an additional effect because perturbing the influent also changes the mass conservation reference. In the reduced example, suppose the driver has no influent term. Increasing either influent component by \(\delta\) increases the required total by \(\delta\). The unweighted affine projection then adds \(\delta/2\) to each predicted component. The resulting influent derivative is \(\frac12\begin{bmatrix}1&1\\1&1\end{bmatrix}\), even though the raw regression has zero direct influent sensitivity. This contribution arises entirely from the changing right-hand side of the physical equality.

Define
\begin{equation}
\equationname{Complementary mass conservation projection matrices}
\label{eq:icsor_projection_matrices}
Q_A=A^\top(AA^\top)^{-1}A,
\qquad P_A=I_F-Q_A.
\end{equation}
Then \(c_{aff}=P_Ac_{raw}+Q_Ac_{in}\), and with the fitted coefficients and invariant operator held fixed,
\begin{equation}
\equationname{Operating and influent Jacobians of the affine-projected state}
\label{eq:icsor_affine_jacobians}
J_{aff,u}=P_AJ_{raw,u},
\qquad J_{aff,c}=P_AJ_{raw,c}+Q_A.
\end{equation}
The \(Q_A\) term in the influent derivative accounts for the changing conservation reference.

These expressions satisfy \(AJ_{aff,u}=0\) and \(AJ_{aff,c}=A\). An operating perturbation therefore changes the affine prediction only in directions that preserve the current influent inventories, while an influent perturbation changes those inventories consistently with the perturbed influent state. These identities also provide analytical checks on the sensitivity implementation.

When the same affine branch remains active over a local neighborhood, these Jacobians describe that branch. The complete deployment rule, however, can transition among raw, affine, and linear-program branches. Each branch therefore has its own local sensitivity \citep{Selerio2026}.

\subsubsection{Active constraints in the final projection}

A non-negativity boundary can change the derivative again. For the same reduced example, the affine components are
\begin{equation}
\equationname{Illustrative affine-projected components as functions of an operating input}
\label{eq:icsor_small_affine_branch}
c_{aff,1}=\frac{35}{6}+\frac{5a}{12},
\qquad
c_{aff,2}=\frac{25}{6}-\frac{5a}{12}.
\end{equation}
Their sum is 10. Near \(a=10\), the first component remains positive while the second reaches zero. Immediately below the boundary, the affine state is feasible and has derivative \((5/12,-5/12)^\top\). Above the boundary, the second affine component becomes negative, and the positive-weight absolute projection returns \((10,0)^\top\). The deployed derivative then becomes \((0,0)^\top\). The final response therefore has different local behavior on the two sides of the active constraint.

The same result follows by differentiating the active equalities. In the projected branch, the binding conditions are \(c_1+c_2=10\) and \(c_2=0\). Differentiation gives
\[
\frac{\partial c_1}{\partial a}+\frac{\partial c_2}{\partial a}=0,
\qquad
\frac{\partial c_2}{\partial a}=0,
\]
which implies a zero derivative for both components. The general active-set calculation extends this idea to all binding conservation and adjustment constraints.

The linear program can activate a zero-concentration bound or one of the absolute-adjustment inequalities. These binding constraints define the active set. If the optimal solution is unique and its nondegenerate basis remains unchanged over a local neighborhood, the solution varies affinely with the relevant right-hand sides. A local branch-specific sensitivity can then be obtained directly from that fixed basis.

At an active-set transition, however, the derivative can change abruptly. If several optimal projected states exist, the returned response can also depend on the rule used to select among them. Local deployment sensitivity is therefore assessed either by differentiating a verified fixed basis or by perturbing admissible inputs and resolving the complete deployment problem. Any branch change is recorded as part of that sensitivity assessment.

This distinction prevents three different interpretive quantities from being conflated. A driver coefficient belongs to \(\widehat B\). A raw-state sensitivity includes the effect of \(\widehat R^{-1}\). A deployed-state sensitivity additionally depends on the mass conservation reference and any active projection constraints. Each describes a different stage of the prediction.

\subsubsection{Following sensitivity through projection}

Interpretation must therefore follow the complete prediction map rather than stop at the fitted polynomial. Coupling transforms the driver into component predictions, the reporting map converts those components into water-quality quantities, and active projection constraints can alter the local derivative of the final state. Physical-unit interpretation should therefore specify whether it concerns a driver coefficient, a raw response derivative, or the derivative of the final projected prediction \citep{Brunton2016,Shen2023}.

A second hypothetical example isolates the effect of projection. Let a dimensionless operating input be \(a\), and define the raw concentrations
\[
c_{raw,1}(a)=4+a,\qquad c_{raw,2}(a)=3+2a.
\]
The raw slopes are one and two. At \(a=2\), the concentrations are 6 and 7, giving a total of 13. Suppose the required total is 10 and the projection minimizes the sum of squared adjustments with equal weights. If the adjustments are \(d_1\) and \(d_2\), conservation requires \(d_1+d_2=-3\). Substituting \(d_2=-3-d_1\) into the objective gives
\[
\begin{aligned}
d_1^2+d_2^2&=d_1^2+(-3-d_1)^2,\\
\frac{\mathrm d}{\mathrm d d_1}\bigl[d_1^2+(-3-d_1)^2\bigr]
&=4d_1+6=0.
\end{aligned}
\]
Thus \(d_1=d_2=-1.5\), giving the mass-conserving state \((4.5,5.5)^\top\).

For general \(a\), the excess total is \((4+a)+(3+2a)-10=3a-3\). Equal-weight affine projection subtracts half of that excess from each coordinate:
\[
\begin{aligned}
c_{aff,1}(a)&=4+a-\frac{3a-3}{2}=5.5-0.5a,\\
c_{aff,2}(a)&=3+2a-\frac{3a-3}{2}=4.5+0.5a.
\end{aligned}
\]
The projected slopes are therefore \(-0.5\) and \(0.5\). The first component changes from a positive raw slope to a negative affine slope, while the two projected slopes sum to zero as required by the fixed total inventory.

Non-negativity introduces a further change. The first affine component reaches zero at \(a=11\). For \(a>11\), the affine state lies outside the non-negative conservation segment, and the nearest jointly feasible state is \((0,10)^\top\). Both local slopes are then zero while that boundary remains active. A single operating input can therefore produce three different local responses: the raw polynomial slope, the affine mass-conserving slope, and the final boundary-constrained slope.

Figure~\ref{fig:assessment_branch_sensitivity} shows this progression for the first component.

\begin{figure}[htbp]
\centering
\includegraphics[width=\textwidth]{ch7_concept_branch_sensitivity.pdf}
\caption{Hypothetical first-component response and sensitivity before and after mass conservation and non-negativity enforcement. Open circles mark the two one-sided checked slopes at the active-boundary transition.}
\label{fig:assessment_branch_sensitivity}
\end{figure}

Before the concentration boundary becomes active, the affine and checked responses coincide. Their negative first-component slope differs from the positive raw slope because mass conservation redistributes the operating response across the two components. After the boundary activates, the checked concentration remains at zero and its local slope becomes zero. The transition is therefore characterized by distinct one-sided sensitivities. This example summarizes why interpretation must follow the complete deployed calculation.

\subsection{Assessment design for the structured regression}

\subsubsection{Common inputs and comparator models}

The structured-regression assessment uses a design collection of 10,000 accepted steady-state reactor cases. Every predictor receives the same 22 physical inputs and targets the same 20 effluent components. COD, TN, TP, and TSS are subsequently calculated through the common composition matrix. The input domain, reaction model, simulation bounds, and acceptance criteria remain those of the declared reactor benchmark. The comparison therefore changes the statistical predictor while keeping the physical process and prediction task fixed.

The nine unconstrained comparators are XGBoost, LightGBM, CatBoost, AdaBoost, random forest, support vector regression, \(k\)-nearest neighbors, partial least squares, and a fully connected multilayer perceptron. Together, they represent boosted and randomized tree ensembles, kernel regression, local neighbor averaging, latent linear regression, and nonlinear neural regression. Their methodological foundations include the ensemble construction of \citet{Breiman2001}, the boosting formulations of \citet{Drucker1997}, \citet{ChenGuestrin2016}, \citet{Ke2017}, and \citet{Prokhorenkova2018}, and the latent regression framework of \citet{Wold2001}. These comparators do not use an invariant penalty, auxiliary non-negative fitted-state matrix, or learned component-coupling equation during fitting. The latent components of partial least squares are statistical combinations of the inputs rather than physical process components. The supervised target nevertheless remains the same 20-component effluent state for all methods.

The neural comparator uses hidden layers of 128, 64, and 32 units with logistic activation. Dense connections map the 22 inputs through these layers to the 20 component outputs. Its flexible internal representation is therefore compared with the explicit polynomial structure of ICSOR without changing the prediction target or reporting map.

\subsubsection{Training partitions and hyperparameter selection}

The primary assessment uses a row-wise 80\% training and 20\% held-out split with random seed 42. From the 10,000-case design, 8000 observations form the training pool and 2000 form the final held-out set. Hyperparameter selection uses only the training pool. Half of that pool forms a 4000-row tuning subset, which is further split into 3200 tuning-training and 800 tuning-validation observations.

Each predictor receives 100 hyperparameter trials under a seeded tree-structured Parzen estimator with seed 42. No trial pruning or wall-clock timeout is imposed. \citet{Bergstra2011} develop this search strategy, while \citet{Akiba2019} provide the associated hyperparameter optimization framework. The selection criterion is mean squared error in the shared physical component space on the internal validation set. Once selected, the settings are held fixed for the primary assessment and repeated sample-size study. The final primary model is then fitted using all 8000 training observations.

For the unconstrained comparators, input and output standardization parameters are estimated from the applicable training partition. Predictions are transformed back to physical component coordinates before invariant checking, reporting, or error assessment. ICSOR remains in physical coordinates throughout so that its coefficients, invariant operator, and deployment constraints share the same component convention. Internal-validation transformations are estimated only from the tuning-training rows. Thus, means, standard deviations, and hyperparameters are determined without access to the held-out test responses. Projection weights are likewise fixed before held-out assessment.

Table~\ref{tab:icsor_comparator_settings} lists the specified comparator configurations. These settings describe the fitted models and should not be interpreted as performance results.

\begin{table}[htbp]
\centering\small
\caption{Specified configurations of the nine unconstrained comparators.}\label{tab:icsor_comparator_settings}
\begin{tabular}{p{0.19\textwidth}p{0.74\textwidth}}
\toprule
Comparator & Configuration\\
\midrule
XGBoost & 309 trees, learning rate 0.125, maximum depth 3, minimum child weight 6.84, row fraction 0.915, feature fraction 0.828, absolute penalty \(1.18\times10^{-3}\), and squared penalty 0.382.\\
LightGBM & 302 trees, learning rate 0.0809, 16 leaves, maximum depth 8, minimum leaf count 14, row fraction 0.845, feature fraction 0.869, absolute penalty \(3.17\times10^{-3}\), and squared penalty \(3.36\times10^{-4}\).\\
CatBoost & 236 boosting iterations, learning rate 0.172, depth 4, leaf squared penalty 0.0266, bagging temperature 0.0277, and random strength 0.0657.\\
AdaBoost & 236 estimators, learning rate 0.999, and squared regression loss.\\
Random forest & 390 trees, maximum depth 14, minimum split count 3, minimum leaf count 1, feature fraction 0.500, and no bootstrap resampling.\\
Support vector regression & Polynomial kernel of degree 4, regularization parameter 0.118, insensitive-error width 0.0458, constant offset 0.889, and solver tolerance \(8.96\times10^{-4}\). The kernel scale is the reciprocal of the product of input dimension and training input variance.\\
\(k\)-nearest neighbors & 15 neighbors, inverse-distance weighting, Manhattan distance, and a k-dimensional search tree with leaf size 40.\\
Partial least squares & 6 latent components, maximum 887 fitting iterations, and tolerance \(9.91\times10^{-5}\).\\
Multilayer perceptron & Hidden widths \((128,64,32)\), logistic activation, adaptive moment estimation, squared penalty \(6.33\times10^{-4}\), batch size 64, initial learning rate \(1.08\times10^{-3}\), maximum 500 iterations, and tolerance \(1.63\times10^{-5}\). Training rows are not shuffled, and early stopping is disabled.\\
\bottomrule
\end{tabular}
\end{table}

\subsubsection{Scoring states and physical diagnostics}

The scored ICSOR prediction is its final checked component state. For comparator accuracy assessment, the raw comparator prediction receives the same affine mass-conservation projection in Equation~\eqref{eq:icsor_affine_solution}, after which the common reporting map is applied. The accuracy comparison therefore evaluates the checked structured prediction against comparator predictions aligned to the same invariant equalities. Comparator accuracy scoring uses the affine state, whereas ICSOR accuracy scoring uses the complete deployed state, including the constrained non-negativity projection when required.

The numerical assessment uses an operator \(A_{num}\) obtained by rounding the null-space basis, removing negligible coefficients, and dividing each row by its first nonzero coefficient. Fitting, affine projection, and constrained deployment use this numerical operator in place of the ideal operator \(A\). Reported balance violations therefore refer specifically to these adopted numerical relations and the stated scaling.

Physical diagnostics use the raw comparator predictions rather than their affine-aligned accuracy states. ICSOR diagnostics distinguish the raw coupled prediction, affine projection, and final deployed state. This separation makes it possible to identify which stage produces physical compliance. The empirical balance diagnostic is \(v_{mc}^{num}(c)=\|A_{num}(c-c_{in})\|_2\), while non-negativity is assessed using \(v_{nn}\) from Equation~\eqref{eq:icsor_deployment_diagnostics}. For either diagnostic,
\begin{equation}
\equationname{Percentage frequency of a physical-constraint violation}
\label{eq:icsor_violation_frequency}
f_v=\frac{100}{n}\sum_{i=1}^{n}
\mathbf{1}\{v(c_i)>10^{-10}\}.
\end{equation}
Violation frequencies, residual magnitudes, adjustment distances, and deployment-branch use answer different questions about the same prediction procedure. They are therefore reported alongside, rather than substituted for, predictive error.

\subsubsection{Composite error measures}

The five reported error measures summarize different aspects of predictive disagreement. Consider two hypothetical observations with reference values \((2,4)\) and predictions \((3,2)\). The residuals are \((1,-2)\), their squares are \((1,4)\), and their absolute values are \((1,2)\). The resulting mean squared error is 2.5, root mean squared error is approximately 1.5811, and mean absolute error is 1.5.

The corresponding absolute relative errors are \(1/2\) and \(2/4\), giving a mean relative absolute error of 0.5. The reference mean is 3, with total squared variation \((2-3)^2+(4-3)^2=2\). The prediction residual squared sum is 5, so \(R^2=1-5/2=-1.5\). The negative value indicates that the hypothetical prediction has greater squared error than using the reference mean for both observations. These calculations illustrate why the metrics provide complementary information.

For reported quantity \(q\in\{1,2,3,4\}\), let \(e_{iq}=\widehat y_{iq}-y_{iq}\). The squared, root, and absolute measures are
\begin{equation}
\equationname{Per-composite squared, root, and absolute error measures}
\label{eq:icsor_composite_error_metrics}
E_{2,q}=\frac{1}{n}\sum_i e_{iq}^{2},
\qquad E_{root,q}=\sqrt{E_{2,q}},
\qquad E_{1,q}=\frac{1}{n}\sum_i|e_{iq}|.
\end{equation}
For strictly nonzero reference values,
\begin{equation}
\equationname{Mean relative absolute error of a reported composite}
\label{eq:icsor_relative_error}
E_{rel,q}=\frac{1}{n}\sum_i\frac{|e_{iq}|}{|y_{iq}|}.
\end{equation}
The quantity is reported as a ratio; multiplying by 100 gives a percentage. A zero reference value makes the expression undefined, while a very small reference can make the relative error disproportionately large. Any convention for zero-reference cases must therefore be defined before predictors are compared.

The target-specific coefficient of determination is
\begin{equation}
\equationname{Coefficient of determination for a reported composite}
\label{eq:icsor_coefficient_determination}
R_q^2=1-
\frac{\sum_i e_{iq}^{2}}
{\sum_i(y_{iq}-\overline y_q)^2}.
\end{equation}
A target with no reference variation makes the expression undefined. Negative values are possible when the prediction error exceeds that of the reference-mean predictor.

The aggregate root error pools residuals across the four reported quantities:
\begin{equation}
\equationname{Root mean squared error pooled across four reported composites}
\label{eq:icsor_aggregate_root_error}
E_{root}=\left(\frac{1}{4n}\sum_{q=1}^{4}\sum_{i=1}^{n}e_{iq}^{2}\right)^{1/2}.
\end{equation}
Aggregate absolute and relative errors use the corresponding means across observations and reported quantities, while aggregate \(R^2\) is the mean of the four target-specific coefficients. Because the composite coordinates retain their own physical units, the pooled error is a benchmark summary under that numerical convention rather than a dimensionless measure. Per-target scores remain necessary to identify which reported quantities drive the aggregate.

Pooling before taking the square root is also distinct from averaging target-specific root errors. If four hypothetical target mean squared errors are 1, 9, 1, and 9, the pooled mean squared error is 5 and the pooled root error is \(\sqrt5\approx2.2361\). The average of the four target root errors is instead \((1+3+1+3)/4=2\). Equation~\eqref{eq:icsor_aggregate_root_error} defines the first of these two summaries.

The primary ordering measure is aggregate root mean squared error. The remaining measures provide complementary perspectives: squared error emphasizes large residuals, absolute error describes average deviation, relative error scales that deviation to the reference magnitude, and \(R^2\) relates squared error to target variation. Predictor rankings are therefore examined separately under each measure.

\subsubsection{Repeated sample-size assessment}

The sample-size assessment uses
\begin{equation}
\equationname{Observation counts for structured-regression sample-size assessment}
\label{eq:icsor_assessment_sizes}
n\in\{500,1450,2400,3350,4300,5250,6200,7150,8100,9050,10000\}.
\end{equation}
At each size, ten subsamples are drawn without replacement from the design collection using deterministic seeds derived from 42. Each subsample receives an 80\% training and 20\% held-out split, and all ten predictors use the same sampled rows and partitions. The complete design therefore contains 110 fits per predictor.

The selected hyperparameters remain fixed across this assessment. This isolates the effect of data availability under a common model specification rather than allowing each sample size to define a new tuned model. Mean held-out scores and their variability are summarized across the ten runs. These learning curves show how the relative value of model flexibility and structural constraints changes with the simulation budget \citep{Geman1992,VieringLoog2023}.

The learning-curve area can be understood from a reduced example. Suppose an error measure takes values 6, 4, and 3 at sample sizes 100, 200, and 400. The first trapezoid has width 100 and average height 5, giving area 500. The second has width 200 and average height 3.5, giving area 700. The total area is therefore 1200. Dividing by the complete width \(400-100=300\) gives a normalized area of 4. Wider sample-size intervals contribute proportionally more to this average.

For retained sizes \(n_1,\ldots,n_S\), the numerical approximation is
\begin{equation}
\equationname{Trapezoidal approximation of normalized learning-curve area}
\label{eq:icsor_learning_area_trapezoids}
\frac{1}{n_{max}-n_{min}}
\sum_{s=1}^{S-1}(n_{s+1}-n_s)
\frac{\overline M_{m,j}(n_s)+\overline M_{m,j}(n_{s+1})}{2}.
\end{equation}
The corresponding continuous definition is
\begin{equation}
\equationname{Integral definition of normalized learning-curve area}
\label{eq:icsor_learning_area}
L_{m,j}=\frac{1}{n_{max}-n_{min}}
\int_{n_{min}}^{n_{max}}\overline M_{m,j}(n)\,\mathrm{d}n.
\end{equation}
Here \(\overline M_{m,j}(n)\) is the mean held-out value of measure \(j\) for predictor \(m\) at sample size \(n\). Dividing by the sample-size interval makes the result an average curve height and preserves the units of the original measure. \citet{Cawley2011} and \citet{Guyon2011} provide related trajectory summaries in active learning, while the present study varies the number of randomly sampled training observations.

Ranks provide a separate view of relative ordering. If four hypothetical errors are 3.0, 3.0, 4.5, and 6.0, their dense ranks are 1, 1, 2, and 3. The tie receives the same rank, and the next distinct score receives the next consecutive rank. The numerical distance between 4.5 and 6.0 does not affect the difference between their ranks.

For each measure \(j\), reported quantity \(q\), and retained sample size \(n_s\), predictors receive a dense rank \(r_{m,j,q}(n_s)\) based on their mean held-out score. Lower squared, root, absolute, and relative errors receive better ranks, whereas larger \(R^2\) values receive better ranks. With \(S=11\) sizes and five measures,
\begin{equation}
\equationname{Mean dense ranks over quantities, measures, and sample sizes}
\label{eq:icsor_dense_rank}
\overline r_{m,j}=\frac{1}{4S}\sum_{q=1}^{4}\sum_{s=1}^{S}r_{m,j,q}(n_s),
\qquad
\overline r_m=\frac{1}{5}\sum_{j=1}^{5}\overline r_{m,j}.
\end{equation}
These summaries describe relative performance over the declared comparisons. \citet{Demsar2006} develops broader rank-based model comparisons, while \citet{Shevchenko2024} emphasize the importance of benchmarking variability. Absolute error and mean rank are therefore reported together so that ordering is not detached from error magnitude.

At the largest sample size, the generalization gap is
\begin{equation}
\equationname{Held-out minus training root-error gap}
\label{eq:icsor_generalization_gap}
\Delta E_{root,m}=E_{root,m}^{held\text{-}out}-E_{root,m}^{training}.
\end{equation}
A positive value indicates greater held-out error. The gap does not by itself measure predictive quality: a model can have a small gap because it performs similarly well on training and held-out data, or because it performs similarly poorly on both. For example, training and held-out errors of 0.5 and 0.8 give a gap of 0.3, while errors of 5.0 and 5.1 give a smaller gap of 0.1 despite much worse prediction. The gap is therefore interpreted together with the absolute error \citep{Hastie2009}.

Training time is recorded for every repeated fit and summarized across runs. Deployment effort is assessed separately through prediction time, branch use, and constrained-solve effort. This separation distinguishes the cost of constructing the approximation from the cost of returning a physically checked state.

\subsection{Interpretability assessment}

The interpretability assessment examines the six driver blocks, the coupling matrix, and the corresponding raw component and composite mappings. It asks what the fitted structure reveals about the response within the declared process domain. \citet{Brunton2016} connect interpretable nonlinear terms to the broader problem of recovering mathematical structure from data, while \citet{Machado2014} show how operational and influent uncertainty complicate interpretation in activated sludge systems.

The driver and coupling parameters can trade off in the simultaneous relation \(Rc=B\phi\). Different joint stationary solutions can therefore produce similar component predictions using different coefficient allocations \citep{Selerio2026}. The imposed symmetry removes one source of redundancy in the ordered quadratic terms, but it does not establish unique causal parameter identification. Coefficient signs and magnitudes are therefore interpreted together with the complete fitted response rather than in isolation.

Prediction error and physical diagnostics are also separated across the raw regression, affine projection, and final deployed state. This distinction makes clear what the statistical model learns and what the projection subsequently imposes \citep{Karniadakis2021,KircherVotsmeier2025}. Projection magnitude is examined because a final state can satisfy the adopted constraints while remaining far from the raw prediction. Likewise, fitting effort and deployment effort are recorded separately because they represent different computational costs.

The assessment is therefore limited to steady-state component prediction within the stated input domain and control volume. The raw regression provides an inspectable statistical approximation. Checked deployment then imposes the adopted numerical invariants and non-negativity requirements on the returned state.

\section{Results}
\label{sec:icsor_results}

\subsection{Accuracy on the fixed held-out partition}

Table~\ref{tab:icsor_fixed_results} summarizes predictive performance on the 2000 fixed held-out cases. ICSOR is evaluated using its final checked state, while the comparator predictions are evaluated after affine alignment to the adopted mass conservation relations. The multilayer perceptron produced the lowest pooled root mean squared error at 4.38, followed by LightGBM at 5.30. Support vector regression and ICSOR were closely spaced at 5.97 and 5.98, respectively. Under the pooled native composite-coordinate convention of Equation~\eqref{eq:icsor_aggregate_root_error}, ICSOR therefore provided competitive but not leading full-data accuracy.

\begin{table}[htbp]
\centering\small
\caption{Fixed-split held-out composite performance. ICSOR is scored after checked deployment. Comparators are scored after affine alignment without a non-negativity constraint. Relative absolute error is a ratio rather than a percentage.}
\label{tab:icsor_fixed_results}
\begin{tabular}{p{4.7cm}rrrr}
\toprule
Predictor & $E_{root}$ & $E_1$ & Mean $R^2$ & $E_{rel}$\\
\midrule
Multilayer perceptron & 4.38 & 2.55 & 0.9915 & 0.0123\\
LightGBM & 5.30 & 3.21 & 0.9806 & 0.0203\\
Support vector regression & 5.97 & 3.49 & 0.9713 & 0.0244\\
ICSOR & 5.98 & 3.70 & 0.9700 & 0.0256\\
XGBoost & 6.12 & 3.81 & 0.9733 & 0.0244\\
CatBoost & 6.45 & 4.07 & 0.9727 & 0.0250\\
AdaBoost & 21.27 & 13.15 & 0.8643 & 0.0653\\
$k$-nearest neighbors & 23.61 & 13.87 & 0.8535 & 0.0557\\
Partial least squares & 28.24 & 17.05 & 0.7808 & 0.0681\\
Random Forest & 29.00 & 16.93 & 0.7934 & 0.0632\\
\bottomrule
\end{tabular}
\end{table}

Relative to the two leading predictors, ICSOR's pooled root error was 36.5\% higher than that of the multilayer perceptron and 12.8\% higher than that of LightGBM. It nevertheless remained slightly below XGBoost and CatBoost on the same measure. Its mean absolute relative error of 0.0256 corresponds to 2.56\%. The fixed-split comparison therefore quantifies the predictive cost associated with using the more structured model under the largest-data setting.

Figure~\ref{fig:icsor_fixed_results} shows the same comparison for pooled root error, pooled absolute error, and mean composite \(R^2\). ICSOR ranked fourth by pooled root error on this partition, reinforcing the conclusion that its principal advantage at the largest design size is not minimum error alone.

\begin{figure}[htbp]
\centering
\includegraphics[width=\textwidth]{results/ch5/q01_fixed_split_accuracy.png}
\caption{Fixed-split composite accuracy on 2000 held-out cases. Teal bars identify ICSOR after checked deployment. Gray bars identify comparators after affine alignment. Pooled root and absolute errors retain the native composite-coordinate convention. The right panel averages four composite-specific $R^2$ values.}
\label{fig:icsor_fixed_results}
\end{figure}

The aggregate score also hides differences among the individual treatment quantities. Table~\ref{tab:icsor_target_results} separates COD, TN, TP, and TSS. ICSOR's COD root error was 7.77, only 0.18 above LightGBM and lower than that of the other non-neural predictors. Its TN and TSS errors were 3.55 and 8.36, respectively. The multilayer perceptron produced the lowest COD, TN, and TSS errors. All methods had the same TP error of 0.0361 after the specified alignment and reporting procedure, so TP did not distinguish the predictors in this comparison.

\begin{table}[htbp]
\centering\small
\caption{Fixed-split root mean squared error for each composite. COD, TN, TP, and TSS use g COD m$^{-3}$, g N m$^{-3}$, g P m$^{-3}$, and g TSS m$^{-3}$, respectively.}
\label{tab:icsor_target_results}
\begin{tabular}{p{4.7cm}rrrr}
\toprule
Predictor & COD & TN & TP & TSS\\
\midrule
Multilayer perceptron & 6.48 & 1.58 & 0.0361 & 5.66\\
LightGBM & 7.59 & 2.77 & 0.0361 & 6.85\\
Support vector regression & 8.87 & 3.50 & 0.0361 & 7.20\\
ICSOR & 7.77 & 3.55 & 0.0361 & 8.36\\
XGBoost & 8.78 & 3.27 & 0.0361 & 7.88\\
CatBoost & 9.13 & 3.23 & 0.0361 & 8.53\\
AdaBoost & 38.30 & 5.58 & 0.0361 & 17.68\\
$k$-nearest neighbors & 38.24 & 4.23 & 0.0361 & 27.37\\
Partial least squares & 43.22 & 5.09 & 0.0361 & 36.00\\
Random Forest & 45.87 & 4.00 & 0.0361 & 35.26\\
\bottomrule
\end{tabular}
\end{table}

Figure~\ref{fig:icsor_composite_results} makes the scale difference among these quantities visible. COD and TSS contributed much larger numerical errors to the pooled score than TP. The separate target panels are therefore necessary for interpreting what the aggregate comparison represents.

\begin{figure}[htbp]
\centering
\includegraphics[width=\textwidth]{results/ch5/q02_composite_accuracy.png}
\caption{Fixed-split root mean squared error for COD, TN, TP, and TSS. Teal identifies ICSOR and gray identifies the comparators. Each panel has its own constituent unit and horizontal scale. All methods have the same reported TP error under the specified alignment and reporting procedure.}
\label{fig:icsor_composite_results}
\end{figure}

\subsection{Sample efficiency and generalization}

The repeated sample-size assessment produced a different picture from the fixed large-data comparison. At the smallest design size of 500 cases, ICSOR had the lowest mean held-out root error at 10.79. XGBoost and LightGBM followed at 11.26 and 11.45, while the multilayer perceptron ranked fifth at 13.26. At the largest design size, however, the ordering changed: the multilayer perceptron reached 4.62, LightGBM 5.27, support vector regression 5.84, and ICSOR 5.94. The leading model therefore depended on the available simulation budget rather than remaining fixed across the learning curve.

\begin{table}[htbp]
\centering\small
\caption{Repeated-split composite assessment over eleven design sizes and ten runs per size. Final error and the held-out minus training gap refer to $n=10{,}000$. The learning-curve area is normalized by the design-size interval. Lower error and curve-area values indicate better performance. The gap measures the separation between training and held-out error.}
\label{tab:icsor_scaling_results}
\begin{tabular}{p{4.3cm}rrrr}
\toprule
Predictor & \begin{tabular}{r}Final\\$E_{root}$\end{tabular} & \begin{tabular}{r}Normalized\\curve area\end{tabular} & Mean rank & $\Delta E_{root}$\\
\midrule
Multilayer perceptron & 4.62 & 6.75 & 2.5 & 0.65\\
LightGBM & 5.27 & 6.35 & 2.2 & 1.99\\
Support vector regression & 5.84 & 6.90 & 4.0 & 1.19\\
ICSOR & 5.94 & 6.33 & 4.3 & 0.22\\
XGBoost & 6.10 & 7.01 & 4.3 & 1.44\\
CatBoost & 6.26 & 7.60 & 5.0 & 0.69\\
AdaBoost & 21.42 & 20.29 & 8.2 & 0.23\\
$k$-nearest neighbors & 23.76 & 25.77 & 7.0 & 23.66\\
Partial least squares & 27.96 & 30.17 & 9.4 & 0.38\\
Random Forest & 28.78 & 30.31 & 8.2 & 18.71\\
\bottomrule
\end{tabular}
\end{table}

Across the complete range of design sizes, ICSOR had the lowest normalized root-error curve area at 6.33. LightGBM was nearly identical at 6.35 and achieved the best mean dense rank of 2.2. The multilayer perceptron had the second-best mean rank at 2.5, while ICSOR and XGBoost both had a mean rank of 4.3. These summaries emphasize different aspects of performance: the curve area favors sustained low root error across the data range, whereas the rank aggregates relative ordering across targets, metrics, and sizes. ICSOR therefore led one trajectory summary without leading all comparative summaries.

Its training-to-held-out behavior also differed from that of the more flexible learners. At the largest design size, ICSOR's held-out minus training root-error gap was 0.22, smaller than the corresponding gaps for the multilayer perceptron, LightGBM, XGBoost, and support vector regression. Partial least squares also had a small gap of 0.38, but its held-out root error was 27.96. A small gap alone therefore did not imply useful prediction. ICSOR combined a small gap with substantially lower absolute error than partial least squares, while the most accurate flexible learners still achieved lower final prediction error.

Figure~\ref{fig:icsor_efficiency_results} separates the normalized curve area, mean dense rank, and final generalization gap. The figure reinforces that no single summary dominates the interpretation: ICSOR led the root-error area, LightGBM led the mean-rank comparison, and the absolute held-out error remained necessary for interpreting the generalization gap.

\begin{figure}[htbp]
\centering
\includegraphics[width=\textwidth]{results/ch5/q03_sample_efficiency.png}
\caption{Normalized root-error curve area, mean dense rank, and held-out minus training root-error gap. Curve areas and ranks summarize eleven design sizes with ten repeated splits per size. The gap refers to the largest design size. Teal identifies ICSOR. The gap is reported alongside the held-out error.}
\label{fig:icsor_efficiency_results}
\end{figure}

The change in model ordering is especially clear at the two endpoints shown in Figure~\ref{fig:icsor_budget_results}. ICSOR led the four selected predictors when only 500 cases were available, whereas the multilayer perceptron had the lowest error at 10,000 cases. The result therefore supports a conditional interpretation of the structured model: its advantage is strongest when simulation data are limited, while flexible nonlinear learners benefit more from the larger design.

\begin{figure}[htbp]
\centering
\includegraphics[width=\textwidth]{results/ch5/q04_data_budget_comparison.png}
\caption{Mean held-out pooled composite root error at design sizes 500 and 10,000 for ICSOR, XGBoost, LightGBM, and the multilayer perceptron. Values summarize ten repeated splits. The panels show repeated-split endpoint means at the two specified design sizes.}
\label{fig:icsor_budget_results}
\end{figure}

\subsection{Physical checks and deployment branches}

The physical diagnostics distinguish statistical prediction from the properties of the returned component state. Table~\ref{tab:icsor_physical_results} compares the final ICSOR output with the raw outputs of the nine comparators. Every final ICSOR state satisfied both the adopted numerical balances and non-negativity tolerance. By contrast, every raw comparator state violated the adopted balance tolerance. Several comparators---AdaBoost, random forest, and \(k\)-nearest neighbors---produced no negative raw states, but this did not make them mass conserving. For the remaining comparators, raw non-negativity violation rates ranged from 46.00\% to 70.10\%.

\begin{table}[htbp]
\centering\small
\caption{Component-state violation frequencies on the 2000 held-out cases. ICSOR uses its final checked state. Comparators use their raw state, not the affine-aligned state used for accuracy scoring. Balance violations refer to $A_{num}$ and the declared $10^{-10}$ tolerance.}
\label{tab:icsor_physical_results}
\begin{tabular}{p{5.2cm}rr}
\toprule
Predictor & Balance violations (\%) & Negative states (\%)\\
\midrule
Multilayer perceptron & 100.0 & 58.20\\
LightGBM & 100.0 & 46.00\\
Support vector regression & 100.0 & 68.65\\
ICSOR & 0.0 & 0.00\\
XGBoost & 100.0 & 64.95\\
CatBoost & 100.0 & 70.10\\
AdaBoost & 100.0 & 0.00\\
$k$-nearest neighbors & 100.0 & 0.00\\
Partial least squares & 100.0 & 48.70\\
Random Forest & 100.0 & 0.00\\
\bottomrule
\end{tabular}
\end{table}

The stage-wise ICSOR diagnostics show how this final compliance was obtained. None of the raw coupled predictions passed both physical checks because all violated the adopted balance tolerance, although 21.7\% were already non-negative. Affine projection restored the balance equalities in every case and simultaneously produced a non-negative state for 436 predictions, or 21.8\% of the held-out set. These states could be returned immediately. The remaining 1564 cases, or 78.2\%, required the weighted linear projection to enforce non-negativity while preserving the balances. That solve required a mean of 23.5 iterations and a maximum of 40. Every resulting final state passed both checks.

Figure~\ref{fig:icsor_deployment_results} separates the raw comparator diagnostics from the internal stages of ICSOR deployment. This separation is essential because the final zero-violation result was not a property of the raw regression. It was produced by the projection procedure. The affine stage supplied balance compliance, while the constrained linear projection supplied the additional correction needed for most cases to satisfy non-negativity.

\begin{figure}[htbp]
\centering
\includegraphics[width=\textwidth]{results/ch5/q05_physical_deployment.png}
\caption{Physical diagnostics and ICSOR deployment on 2000 held-out cases. The left panel compares final ICSOR states with raw comparator states. The middle panel separates raw coupled, affine, and final ICSOR compliance. The right panel gives the numbers returned after affine projection and after constrained linear projection. Balance checks use the adopted numerical operator $A_{num}$.}
\label{fig:icsor_deployment_results}
\end{figure}

The branch counts therefore identify the source of physical admissibility. Fitting guidance alone did not make the raw state compliant. The affine projection restored the adopted invariant relations, and the constrained projection handled the remaining concentration violations. The complete checked deployment procedure, rather than the regression by itself, produced the zero final violation rates.

\subsection{Computational effort}

The structured fitting procedure required more computation than the fastest comparators, but its relative cost changed with the design size. At 500 cases, mean ICSOR fitting time was 21.5 s, compared with 4.72 s for LightGBM, 3.16 s for XGBoost, and 4.99 s for the multilayer perceptron. At 10,000 cases, ICSOR required 94.4 s. LightGBM and the other fast tree ensembles remained much faster, while partial least squares and \(k\)-nearest neighbors required only 0.228 and 2.70 s, respectively.

ICSOR was nevertheless not the slowest model at the largest design size. Its mean fitting time remained below AdaBoost at 117.7 s, support vector regression at 136.8 s, and the multilayer perceptron at 161.6 s. The result places the structured fit between the fastest conventional regressors and several computationally heavier nonlinear alternatives.

Fitting time is only one part of the computational requirement. Most fixed-split ICSOR predictions also required constrained deployment: the linear projection branch was activated in 78.2\% of cases. The reported iteration counts therefore represent additional repeated-use effort beyond the raw regression evaluation. Separating fitting from deployment makes clear that model preparation and physically checked prediction impose different computational costs.

Figure~\ref{fig:icsor_fitting_results} shows the fitting-time comparison on logarithmic axes. The smaller-design panel focuses on four leading predictors, while the larger-design panel includes additional alternatives. The figure therefore describes model-construction effort only; projection and checked deployment are accounted for separately.

\begin{figure}[htbp]
\centering
\includegraphics[width=\textwidth]{results/ch5/q06_fitting_effort.png}
\caption{Mean fitting times for the specified predictors at design sizes 500 and 10,000. Teal identifies ICSOR. Each panel uses a logarithmic seconds axis. The values concern repeated model fits and exclude deployment calculations.}
\label{fig:icsor_fitting_results}
\end{figure}

\subsection{Inspectable coefficient structure}

The fitted coefficient structure shows where ICSOR remained selective and where it relied on dense dependence among component channels. Table~\ref{tab:icsor_coefficient_results} reports the COD raw-composite polynomial coefficients together with the shared coupling matrix. The polynomial representation contains 276 unique candidate terms. Adding the 380 off-diagonal coupling entries gives 656 entries in the displayed structure. Under the specified post-fitting thresholds, 469 entries were retained, corresponding to 71.5\% of the combined display.

\begin{table}[htbp]
\centering\small
\caption{Coefficients exceeding the reporting thresholds in one fitted model. Polynomial blocks describe the COD row of $B_{\mathrm{raw,comp}}$. The coupling matrix is shared across component outputs and uses its own threshold. Symmetric quadratic blocks count unique upper-triangular terms.}
\label{tab:icsor_coefficient_results}
\begin{tabular}{p{6.2cm}rrr}
\toprule
Coefficient block & Candidates & Retained & Retained (\%)\\
\midrule
Baseline $b$ & 1 & 1 & 100.0\\
Operating terms $W_u$ & 2 & 2 & 100.0\\
Influent terms $W_{in}$ & 20 & 19 & 95.0\\
Operating curvature $\Theta_{uu}$ & 3 & 3 & 100.0\\
Operating-load interactions $\Theta_{uc}$ & 40 & 29 & 72.5\\
Influent curvature $\Theta_{cc}$ & 210 & 44 & 21.0\\
Shared off-diagonal coupling $\Gamma$ & 380 & 371 & 97.6\\
Combined display & 656 & 469 & 71.5\\
\bottomrule
\end{tabular}
\end{table}

The contrast between influent curvature and component coupling was pronounced. Only 44 of 210 unique influent-curvature terms exceeded the display threshold, whereas 371 of 380 coupling entries did. The fitted representation was therefore selective in its nonlinear influent curvature but dense in the dependence learned among output components. This pattern is important for interpretation because it shows that the visible polynomial structure was not uniformly sparse.

The largest absolute COD influent-curvature coefficient was the dissolved-oxygen self term at \(-5.8097\). The two mirrored dissolved-oxygen--dinitrogen entries were each \(-1.7411\), giving a combined unordered interaction coefficient of \(-3.4822\). Other prominent terms linked dissolved oxygen with metal phosphate, ammonia-oxidizing biomass, soluble phosphate, nitrite-oxidizing biomass, and nitrite. These coefficients belong to the raw COD composite mapping after both component coupling and reporting have been applied.

Figure~\ref{fig:icsor_structure_results} shows the retained counts and selected oxygen-related curvature terms. The plotted interaction values combine the two mirrored cells, consistent with the symmetry convention. The figure therefore describes the fitted polynomial in its native input units rather than implying that coefficient magnitudes from different units are directly comparable.

\begin{figure}[htbp]
\centering
\includegraphics[width=\textwidth]{results/ch5/q07_coefficient_structure.png}
\caption{Coefficient structure of one fitted ICSOR model. The left panel separates entries above and below the display thresholds for the COD polynomial blocks and shared coupling. The right panel shows selected COD influent-curvature coefficients with mirrored interaction entries combined. Coefficient units depend on the corresponding input terms. The displayed values describe the raw composite mapping.}
\label{fig:icsor_structure_results}
\end{figure}

\section{Discussion}
\label{sec:icsor_discussion}

\subsection{A structured approximation under limited data}

The results show that the value of the structured model depends strongly on the available data budget. ICSOR produced the lowest mean held-out root error at the smallest design size and the lowest normalized root-error curve area across the full sample-size range. Its explicit second-order driver, invariant-aware fitting objective, and learned component coupling therefore provided a useful approximation when fewer mechanistic simulations were available. This behavior is consistent with the broader role of model structure in hybrid learning discussed by \citet{Shen2023} and with the changing learner orderings across data budgets reviewed by \citet{VieringLoog2023}.

That advantage did not persist as a universal ranking. At the largest design size, the multilayer perceptron produced the lowest composite error, while LightGBM achieved the best mean dense rank across targets, metrics, and sample sizes. ICSOR's normalized curve area of 6.33 was only slightly lower than LightGBM's 6.35, but its final full-data error remained higher than that of the strongest flexible learners. The result therefore does not support choosing ICSOR simply because it is interpretable. Instead, it identifies a specific trade-off: ICSOR offers an explicit equation, favorable performance when simulations are limited, and checked component states, while more flexible learners can achieve lower prediction error when a larger training design is available \citep{Rudin2019,BhosekarIerapetritou2018}.

The generalization results reinforce the same point. ICSOR's training-to-held-out root-error gap was small, but a small gap is useful only when the underlying prediction error is also acceptable. Partial least squares demonstrated the opposite case: it also had a relatively small gap, yet its held-out error remained high. The leading flexible learners achieved lower absolute error despite larger training-to-held-out differences. Generalization gap and predictive accuracy therefore describe complementary properties and should not be interpreted as substitutes for one another.

From an engineering perspective, the sample-size results provide a practical model-selection rule rather than a universal ranking. When generating mechanistic training cases is costly and only a modest design can be afforded, the structured response can make efficient use of those cases while retaining direct inspectability. When simulation data are abundant and minimum predictive error is the dominant requirement, the results favor the more flexible learners evaluated here. The usefulness of an interpretable surrogate therefore depends on the data available and on whether transparency and physical checking are themselves required properties of the engineering workflow.

\subsection{Fitting guidance and final-state enforcement}

The physical results clarify the different roles of fitting guidance and output enforcement. The invariant penalty influenced the auxiliary training states, and the learned coupled system produced the raw prediction for a new input. Neither step, however, guaranteed that the raw deployed state satisfied the adopted physical requirements. Every raw ICSOR prediction in the fixed held-out set required at least the affine mass-conservation projection.

Projection then supplied the guarantees that fitting alone did not provide. The affine stage restored the adopted balances for every held-out case and was sufficient to return a non-negative state for 21.8\% of the predictions. The remaining 78.2\% required the weighted constrained projection, after which every returned state satisfied both the balance and non-negativity checks. The final zero-violation result is therefore attributable to the checked deployment procedure rather than to the soft penalty used during model fitting. This separation is consistent with the distinction between physics-guided learning and explicit output enforcement examined by \citet{Hansen2024} and \citet{KircherVotsmeier2025}.

This distinction is important because a physically checked output is not the same as a statistically superior output. The same deployment procedure that guarantees the adopted balances can alter concentrations and therefore alter prediction error. The numerical comparison used checked ICSOR states and affine-aligned comparator states for accuracy scoring, while the physical comparator diagnostics intentionally used the raw comparator states. Under this convention, ICSOR had a 36.5\% larger pooled root error than the multilayer perceptron, and all methods shared the same TP error of 0.0361 after alignment and reporting. These outcomes should therefore be read together with the deployment rules that produced them.

The result also shows why physical diagnostics need to identify the stage being assessed. A raw prediction, an affine mass-conserving state, and a jointly constrained final state are different mathematical objects. Reporting only the final zero-violation rate would hide how much correction was required to obtain it. Conversely, reporting only raw violations would ignore the state that the model actually returns to the engineer. The stage-wise analysis makes both facts visible: the statistical model frequently required correction, and the complete deployed procedure successfully supplied that correction.

\subsection{What the coefficient structure explains}

The fitted coefficient blocks make the response more inspectable by separating operating dependence, influent dependence, curvature, operating-load interactions, and component coupling. This explicit organization follows the preference for transparent predictive structures emphasized by \citet{Rudin2019}. It also prevents a regression coefficient from being interpreted outside the calculation in which it operates. A coefficient belongs first to the driver, after which coupling and the reporting map can redistribute its effect across component and composite outputs.

The fitted structure was not uniformly sparse. Most unique influent-curvature terms fell below the display threshold, while almost every off-diagonal coupling coefficient exceeded its separate threshold. The displayed model therefore combined selective nonlinear dependence on the influent with dense learned dependence among effluent components. The 469 retained entries out of 656 displayed candidates describe this contrast, but the threshold itself does not alter the fitted model. It only controls which coefficients are made visible in the summary.

The prominent oxygen-related terms provide a useful example of how the structure can be interpreted without claiming causal parameter identification. Oxygen availability is central to nitrification and interacts with other nutrient-removal pathways \citep{StenstromSong1991,Wentzel1992}. The dissolved-oxygen self term and the oxygen--dinitrogen interaction therefore identify fitted response patterns that can be compared with known process behavior. Their signs and magnitudes remain properties of the statistical mapping in its native units, not estimates of kinetic constants.

The sensitivity analysis extends this interpretation beyond coefficient inspection. A coefficient describes one term in the polynomial, but the local response at an operating point combines linear effects, curvature, interactions, component coupling, and physical projection. The analytical raw Jacobians in Equation~\eqref{eq:icsor_raw_jacobians} include the first three of these stages, while Equation~\eqref{eq:icsor_affine_jacobians} shows how mass conservation modifies those derivatives. The frequent activation of the constrained branch means that active-set effects are also relevant for many final predictions.

Figure~\ref{fig:assessment_branch_sensitivity} demonstrates why this distinction matters. In the illustrative response, a positive raw first-component slope became negative after affine mass-conservation projection and then became zero when the non-negativity boundary activated. A reader who inspected only the raw polynomial coefficient would therefore describe a different local response from the one returned by the deployed model. The interpretability contribution of ICSOR lies in making this entire sequence traceable.

\subsection{Engineering implications}

Taken together, the results define the engineering role of the structured surrogate more clearly than any single accuracy score. ICSOR combined competitive composite prediction, favorable behavior when training data were limited, an explicit response structure, and final component states that satisfied the imposed numerical checks. Its second-order terms and coupling matrix allow an engineer to examine how operating and loading conditions enter the fitted response. Checked deployment then ensures that COD, TN, TP, and TSS are calculated from a component state that satisfies the adopted material accounting and non-negativity requirements.

These benefits are accompanied by computational costs. ICSOR required more fitting time than the fastest tree and linear comparators, although at 10,000 cases it remained faster than AdaBoost, support vector regression, and the multilayer perceptron. More importantly, fitting time does not describe the full cost of using the model. Most fixed-split predictions required the constrained deployment solve after the raw regression had already been evaluated. The computational value of the model must therefore be assessed through both preparation and repeated deployment rather than through regression evaluation alone. This distinction follows the broader surrogate-modeling emphasis on separating model-development cost from repeated-use cost \citep{BhosekarIerapetritou2018,Bliek2023}.

The practical value of ICSOR is therefore conditional rather than absolute. It is not the most accurate predictor at the largest data budget, nor is its raw prediction automatically physically admissible. Its value is that the complete procedure exposes the fitted relationships, identifies where coupling and projection alter those relationships, and returns a component state that passes the declared physical checks. When transparency, limited simulation data, and component-level physical accounting are important, those properties provide engineering criteria that are not captured by prediction error alone.

The broader implication is that interpretability and physical consistency should be treated as properties of the complete deployed calculation. An explicit equation does not remain fully interpretable if later coupling and projection are ignored, and a physically constrained output does not become accurate merely because its balances close. The structured framework makes these distinctions explicit. It allows prediction accuracy, model structure, physical enforcement, and computational effort to be examined together without treating any one of them as evidence for the others.

\subsection{Data Availability}
\label{sec:icsor_data_availability}

The simulation data and numerical results supporting this study are publicly available at \url{https://github.com/egbertoselerio1997/icsor}.